% =====================================================================
%  AI-Driven Observability and Root-Cause Analysis
%  Design and Implementation of the CIRES Observability Platform
%
%  Final-Year-Project Report (working draft)
%  Author : Lina Laaraich
%  Institution : Tanger Med / CIRES Technologies — Digital Factory
%  Compiled  : May 2026
% =====================================================================

\documentclass[11pt,a4paper,oneside]{report}

% ---------- packages ----------
\usepackage[utf8]{inputenc}
\usepackage[T1]{fontenc}
\usepackage{lmodern}
\usepackage{amsmath}
\usepackage{amssymb}
\usepackage[english]{babel}
\usepackage[a4paper,margin=2.6cm,headheight=14pt]{geometry}
\usepackage{microtype}
\usepackage{parskip}
\usepackage{enumitem}
\setlist{itemsep=2pt, topsep=4pt}
\usepackage{booktabs}
\usepackage{longtable}
\usepackage{array}
\usepackage{tabularx}
\usepackage{xcolor}
\definecolor{accent}{HTML}{1F6FB0}
\definecolor{accentdark}{HTML}{0B3E73}
\definecolor{boxbg}{HTML}{F4F6FA}
\definecolor{boxrule}{HTML}{C8D2E0}
\definecolor{codebg}{HTML}{F7F7F4}
\definecolor{warnrule}{HTML}{C77B23}
\definecolor{winrule}{HTML}{2A8A4B}
\usepackage{graphicx}
\usepackage{titlesec}
\titleformat{\chapter}[display]
  {\normalfont\bfseries\color{accentdark}}
  {\Large CHAPTER \thechapter}{12pt}{\LARGE}
\titlespacing*{\chapter}{0pt}{-10pt}{14pt}
\usepackage{fancyhdr}
\pagestyle{fancy}
\fancyhf{}
\renewcommand{\headrulewidth}{0.4pt}
\fancyhead[L]{\nouppercase{\leftmark}}
\fancyhead[R]{\thepage}
\fancypagestyle{plain}{\fancyhf{}\fancyfoot[C]{\thepage}\renewcommand{\headrulewidth}{0pt}}
\usepackage{caption}
\captionsetup{font=small,labelfont=bf,labelsep=period}
\usepackage{listings}
\lstset{
  basicstyle=\ttfamily\footnotesize,
  backgroundcolor=\color{codebg},
  frame=single,
  framesep=4pt,
  rulecolor=\color{boxrule},
  breaklines=true,
  showstringspaces=false,
  columns=fullflexible,
  keepspaces=true,
  xleftmargin=8pt,
  xrightmargin=8pt
}
\usepackage{tcolorbox}
\tcbuselibrary{skins,breakable}
\newtcolorbox{notebox}[1][]{enhanced,breakable,
  colback=boxbg,colframe=boxrule,boxrule=0.4pt,arc=2pt,
  left=8pt,right=8pt,top=6pt,bottom=6pt,
  fonttitle=\bfseries\small,#1}
\newtcolorbox{warnbox}[1][]{enhanced,breakable,
  colback=warnrule!6,colframe=warnrule,boxrule=0.5pt,arc=2pt,
  left=8pt,right=8pt,top=6pt,bottom=6pt,
  fonttitle=\bfseries\small\color{warnrule},#1}
\newtcolorbox{winbox}[1][]{enhanced,breakable,
  colback=winrule!5,colframe=winrule,boxrule=0.5pt,arc=2pt,
  left=8pt,right=8pt,top=6pt,bottom=6pt,
  fonttitle=\bfseries\small\color{winrule},#1}
\usepackage[hidelinks,colorlinks=true,
  linkcolor=accent,citecolor=accent,urlcolor=accent]{hyperref}
\usepackage{bookmark}

% ---------- metadata ----------
\newcommand{\reporttitle}{AI-Driven Observability and Root-Cause Analysis}
\newcommand{\reportsubtitle}{Design and Implementation of the CIRES Observability Platform}
\newcommand{\reportauthor}{Lina Laaraich}
\newcommand{\reportorg}{CIRES Technologies / Tanger Med --- Digital Factory Observability Service}
\newcommand{\reportdate}{May 2026}

\hypersetup{
  pdftitle={\reporttitle\ --- \reportsubtitle},
  pdfauthor={\reportauthor},
  pdfsubject={Final-Year Project Report --- AI-driven observability and RCA},
  pdfkeywords={observability, SRE, AIOps, RCA, MCP, Prometheus, Loki, Jaeger, Grafana, LLM, qwen2.5}
}

% ---------- title page ----------
\begin{document}
\begin{titlepage}
\centering
\vspace*{2.5cm}
{\fontsize{12}{14}\selectfont Final-Year Project Report \par}
\vspace{0.6cm}
{\color{accentdark}\rule{0.55\textwidth}{0.6pt}\par}
\vspace{1.0cm}
{\fontsize{24}{30}\selectfont\bfseries \reporttitle \par}
\vspace{0.6cm}
{\fontsize{16}{20}\selectfont\itshape \reportsubtitle \par}
\vspace{0.8cm}
{\color{accentdark}\rule{0.55\textwidth}{0.6pt}\par}
\vspace{2.2cm}
{\large \reportauthor \par}
\vspace{0.4cm}
{\normalsize \reportorg \par}
\vfill
{\large \reportdate \par}
\vspace{0.6cm}
{\small Working draft, prepared as the structural and content basis for the final submitted report.}
\end{titlepage}

% ---------- abstract ----------
\chapter*{Abstract}
\addcontentsline{toc}{chapter}{Abstract}
\markboth{Abstract}{Abstract}

This dissertation reports on the design and implementation of the
\emph{CIRES Observability Platform}, an end-to-end production
observability stack augmented with an AI-driven Root-Cause Analysis
(RCA) layer.  The platform was built between February and May 2026 at
CIRES Technologies (Tanger Med, Digital Factory Observability Service)
across four planned sprints.  Sprints 1--3 were completed during the
project period; Sprint 4 is documented as scoped future work.

The work addresses an operational problem surfaced during the discovery
phase: incidents in the target production environment were being
discovered \emph{after the fact} --- by clients, or during retrospective
log review --- rather than by the monitoring system itself.  The
platform's response is a two-layer design: a conventional
observability foundation (Prometheus, Loki, Jaeger, Grafana, OpenTelemetry)
provisioned by Ansible on AWS, with a FastAPI-based AI triage service
on top that consumes alerts via Grafana webhooks, gathers cross-pillar
context through five Model-Context-Protocol (MCP) bridge servers, and
produces structured root-cause analyses with a local Large Language
Model (\texttt{qwen2.5:7b-instruct} on Ollama).

The implementation introduces several novel mechanisms to keep the
LLM output operationally trustworthy: an exemplar library of fourteen
canonical RCA archetypes injected as structural scaffolding before
the evidence pack; an alert-keyed hallucination blocklist; a
surface-only and hypothesis-menu validator pair; a confidence-clamp
that strips suggested actions when trust signals are weak and emits
read-only \emph{diagnostic steps} instead; a single bounded-agency
retry constrained to one MCP call; and a closed-loop feedback layer
that turns operator overrides into routing signal.  An
\emph{MCP-only data-access invariant} is enforced as a project-wide
rule: every datum the LLM sees must arrive via an MCP bridge, never
via direct in-process imports or direct database/Prometheus queries.

The dissertation presents the architecture, sprint-by-sprint
implementation history, an in-depth case study of the
\texttt{0b215ef3} incident that drove the hallucination-firewall epic,
an evaluation against a 4-axis RCA quality rubric, and a catalogue of
the twenty-three durable design decisions (D1--D23) recorded during
construction.  The final chapter scopes future work, including
iterative agentic gather, curated retrieval-augmented generation
fed from operator feedback, and a labelled chaos-test corpus for
regression-grade quality measurement.

\vspace{1em}
\noindent\textbf{Keywords:} observability, SRE, AIOps, root-cause
analysis, MCP, Prometheus, Loki, Jaeger, Grafana, large language
models, Ollama, qwen2.5, Ansible, Terraform, k3s.

% ---------- TOC ----------
\tableofcontents
\listoffigures
\listoftables

% =====================================================================
\chapter{Introduction}
\label{ch:intro}

\section{Context and motivation}

CIRES Technologies is the digital factory of Tanger Med, the principal
port complex of northern Morocco.  The Digital Factory Observability
Service is responsible for the operational health of an internal
service catalogue that supports port logistics, customer integrations,
and back-office workflows.  To prototype an improved observability
posture without disturbing production systems, the project's work was
deliberately scoped against a \emph{representative web stack} (a
Spring Boot back-end, a Kong API gateway, a MySQL data store, a
React-based front-end) deployed on a small fleet of on-premises
virtual machines --- a scaled-down environment for prototyping the
architecture before any transition into CIRES production; a parallel
migration of the prototype to AWS was already on the medium-term
roadmap.

Internal monitoring at the start of the project was limited to ad-hoc
log inspection and a small number of metric-threshold rules.
Operationally, this produced a structurally reactive posture: outages
were detected by client reports or during retrospective log review,
not by the system itself.  The asymmetry between the speed of failure
(seconds to minutes) and the speed of human detection (hours to days)
defined the operational gap that this project sets out to close.

\section{Problem statement}

\begin{notebox}[title={Problem statement (derived from Sprint 0 field research)}]
The team needs a way to detect production problems early --- before
they reach a client, and without requiring a human to be reading logs
at the right moment.  Detection should also be \emph{pre-investigated}:
each alert should arrive with a candidate root cause, the supporting
evidence, and a recommended remediation, so that the operator's first
move is to validate or override rather than to start an investigation
from scratch.
\end{notebox}

This formulation separates the problem into two distinct halves.  The
first half is conventional: stand up a real observability stack so the
system emits actionable signals.  The second half is novel: layer on
top of that stack an automated reasoning system that turns alerts into
structured, actionable RCAs without on-call human attention.  The
platform implements both halves and integrates them end-to-end.

\section{Aim and objectives}

\textbf{Aim.}  Design, implement and validate a production-shaped
observability platform that proactively detects, investigates and
reports on incidents using an AI triage layer that runs over open-source
infrastructure, with no dependency on hosted SaaS observability
vendors.

\textbf{Objectives.}
\begin{enumerate}
  \item Establish a three-pillar observability foundation (metrics,
        logs, traces) for the target web stack, deployable from clean
        infrastructure via a single Ansible playbook.
  \item Migrate the workload from on-premises VMs onto a
        Terraform-managed AWS environment and run the application
        workload on Kubernetes (k3s) while keeping the monitoring stack
        on Docker Compose for operational simplicity.
  \item Implement an AI triage service that consumes Grafana webhooks
        and produces structured root-cause analyses with verdicts,
        evidence, and suggested remediation actions, persisted to a
        durable history store.
  \item Bridge each telemetry source (Prometheus, Loki, Jaeger,
        triage history, deployment events) behind an MCP server, and
        enforce an MCP-only data-access invariant so that every datum
        the LLM consumes is provenanced through a bridge.
  \item Develop and evaluate quality-control mechanisms that keep
        LLM output operationally trustworthy: exemplar injection, a
        hallucination blocklist, prose-shape validators, a confidence
        clamp, and a closed-loop operator feedback channel.
  \item Document the system end-to-end at a level sufficient for the
        engineering team to operate, extend and audit it after the
        project's conclusion.
\end{enumerate}

\section{Scope and limitations}

The project explicitly scopes \emph{in} the components above and
explicitly scopes \emph{out} the following:

\begin{itemize}
  \item Hosted SaaS observability vendors (Datadog, New Relic,
        Honeycomb).  These were costed and surveyed in Sprint 0; the
        decision to use open-source infrastructure was driven by
        budget, data-sovereignty, and learning objectives.
  \item Production rollout to actual customer-facing systems.  The
        scope is the representative web stack used as a development
        surrogate (a derivative of the
        \texttt{mukundmadhav/react-springboot-mysql} repository),
        deployed on dedicated infrastructure with synthetic load and
        chaos-injected failure.
  \item Multi-region high-availability for the observability stack
        itself.  The platform runs in \texttt{us-east-1} (or local
        on-premises) and is not multi-AZ.
  \item A cloud GPU runner for the LLM layer.  This was originally
        scoped as a Sprint-4 candidate and was dropped on cost grounds
        on 2026-05-19 in favour of a laptop-based runner that proved
        sufficient (see \S\ref{sec:gpu-dropped}).
\end{itemize}

\section{Structure of the dissertation}

The dissertation proceeds as follows.  Chapter~\ref{ch:background}
reviews the technical background: observability practice, SRE,
log-template clustering, large language models in IT operations, and
the Model-Context-Protocol specification.  Chapter~\ref{ch:discovery}
covers the Sprint-0 discovery phase, including interview findings and
the framing that produced the problem statement above.
Chapter~\ref{ch:architecture} presents the platform architecture in
its current shape.  Chapters \ref{ch:sprint1}--\ref{ch:sprint3} cover
the three completed implementation sprints in chronological order.
Chapter~\ref{ch:case-study} provides an in-depth case study of the
\texttt{0b215ef3} incident that drove the hallucination-firewall
epic.  Chapter~\ref{ch:evaluation} evaluates the platform against a
four-axis RCA quality rubric, the chaos-injection harness, and the
recurring static and live audit cycles.  Chapter~\ref{ch:decisions}
catalogues the twenty-three durable design decisions (D1--D23)
recorded during construction and discusses the three most consequential
in depth.  Chapter~\ref{ch:future} scopes Sprint 4 and beyond, and
chapter~\ref{ch:conclusion} concludes.

% =====================================================================
\chapter{Background and Related Work}
\label{ch:background}

\section{The three pillars of observability}

Modern observability practice identifies three primary telemetry
pillars: \emph{metrics}, \emph{logs}, and \emph{traces}.  Metrics are
numeric time-series of system state (CPU utilisation, request rate,
error rate, etc.), typically scraped on a fixed interval and stored in
a purpose-built time-series database such as Prometheus.  Logs are
discrete textual events emitted by services, ideally with structured
fields, aggregated by a log database such as Grafana Loki.  Traces are
distributed records of a single request's journey across services,
carrying a trace identifier that links spans recorded by each hop into
one structure for examination; the OpenTelemetry project provides the
canonical wire format and Jaeger is a widely-used storage backend.

A fully observable system requires all three pillars and the ability
to navigate between them: from an anomalous metric to the logs of the
service whose metric is anomalous, and from those logs to the trace
that carried the originating request.  The CIRES platform implements
this triangulation as a first-class concern; the AI triage service's
context-gathering step explicitly fans out to all three pillars and
to a fourth, Drain3-based log-template clustering signal, before any
LLM reasoning takes place.

\section{SRE practice and the golden signals}

Google's Site Reliability Engineering literature (the
\emph{SRE Book} and \emph{SRE Workbook}) identifies four
\emph{golden signals} that almost any user-facing service should be
measured against: \textbf{latency}, \textbf{traffic},
\textbf{errors}, and \textbf{saturation}.  The book also articulates
two ideas that influenced this project's alert taxonomy:

\begin{itemize}
  \item The \textbf{symptoms vs.\ causes} distinction.  Alerts should
        fire on user-observable symptoms (high latency, elevated error
        rate) rather than on underlying causes (high CPU, high
        memory).  Cause-based alerts produce a flood of false positives
        in healthy systems; symptom-based alerts page the operator
        only when something is actually broken from the user's
        viewpoint.
  \item The \textbf{error budget} framing.  Service-Level Objectives
        (SLOs) define a tolerable error rate; the gap between target
        reliability and 100\,\% reliability is the error budget.  A
        team that is in the red on its error budget should slow down
        feature development and invest in reliability; a team in the
        green has explicit permission to take more risk.
\end{itemize}

The CIRES platform does not at this stage formalise SLOs or error
budgets, but the alert rule set is deliberately structured to lean
symptomatic --- \texttt{HighP95Latency} on the gateway,
\texttt{HighKongP95Latency} per upstream, \texttt{ErrorRate} on the
service tier --- with cause-shaped alerts (CPU, memory, disk)
treated as supplementary diagnostic signal rather than primary paging
material.

\section{Prometheus, Loki, Jaeger, Grafana}

The platform's foundational stack is the de-facto open-source
combination for modern infrastructure observability:

\begin{itemize}
  \item \textbf{Prometheus} for metrics.  Pull-based scraping on a 15-s
        interval.  Spring Boot exposes \texttt{/actuator/prometheus},
        Kong exposes its statistics endpoint, and host-level
        \texttt{node\_exporter} and \texttt{cAdvisor} fill in the gaps.
  \item \textbf{Loki} for logs.  Pull-based ingestion via Promtail.
        Loki stores log streams indexed by labels rather than by
        full-text content, which keeps storage costs low and pairs
        cleanly with Prometheus's label model.
  \item \textbf{Jaeger} for traces.  OpenTelemetry Collector receives
        spans from the OTel Java agent injected into Spring Boot, and
        from the Kong OTel plugin at the gateway layer, and forwards
        them to Jaeger.  Trace identifiers propagate end-to-end.
  \item \textbf{Grafana} for visualisation and alerting.  Grafana
        unified alerting replaced Alertmanager early in the project
        (see decision D6 in chapter~\ref{ch:decisions}); a single
        \texttt{triage-webhook} contact point routes all alert events
        to the AI triage service.
\end{itemize}

\section{Drain3 and log-template clustering}

Real production logs contain a small number of recurring template
strings and a large body of incidental variation (timestamps,
identifiers, parameter values).  The Drain3 algorithm (a tunable
variant of Drain) ingests log lines, replaces the variable parts with
placeholders, and clusters lines by the resulting template.  Once the
catalogue of templates has stabilised, novel templates become a strong
operational signal: a log line whose template has never been seen
before is, by construction, unusual.

The CIRES platform integrates Drain3 as a fourth signal alongside
metrics, logs, and traces.  A background poller reads the most recent
log lines from Loki, feeds them through a per-service template miner,
and emits a synthetic \texttt{Drain3AnomalyDetected} webhook into the
triage service whenever the recent novelty rate exceeds a configurable
threshold.  The webhook payload carries the actual novel template
strings and a sample of the underlying lines, so the LLM downstream
sees not just a count but the evidence itself
(decision~D23 records the per-service split that resolved a class of
false-positive issues introduced by an earlier global miner).

\section{Large language models in IT operations}

The use of large language models for IT-operations tasks is an active
area sometimes labelled \emph{AIOps}.  The state of the practice
varies widely.  Hosted observability vendors have shipped AI-flavoured
features for several years (Datadog Watchdog, New Relic AI,
Honeycomb's BubbleUp, Komodor's incident summarisation), each with a
different scope: anomaly detection, alert summarisation, correlation
across signals, incident reconstruction.

Three risks dominate when an LLM is placed in the operational path:

\begin{enumerate}
  \item \textbf{Hallucination.}  The model may confidently produce
        statements that are not entailed by the evidence pack ---
        invented metric values, services that do not exist, kubectl
        commands for systems that are not deployed on Kubernetes.
  \item \textbf{Surface-only reasoning.}  The model may restate the
        alert in different words (``the alert indicates high latency''
        or ``based on the observed value of 1.2\,s'') without
        identifying a cause.
  \item \textbf{Hypothesis menus.}  Under uncertainty, the model may
        emit a list of possibilities (``this could be a slow query,
        a connection pool issue, or a garbage-collection pause'')
        rather than committing to one and naming the evidence for it.
\end{enumerate}

A central engineering challenge of the project has been the design
of mechanisms that detect and reject each of these failure modes
without rejecting genuinely good RCAs.  Chapter~\ref{ch:sprint3} and
chapter~\ref{ch:case-study} treat these mechanisms in depth.

\section{The Model-Context-Protocol}

The Model-Context-Protocol (MCP) is an open specification, introduced
in 2024 by Anthropic and now adopted across a range of LLM tooling
projects, for connecting LLM applications to external data sources
and tools.  An MCP \emph{server} exposes a small number of tools
(typically read-only) over a JSON-RPC transport; an MCP \emph{client}
discovers tools, invokes them, and feeds the results back into the
model's context.

The CIRES platform uses MCP not in its original ``LLM directly calls
tools'' shape but in a more constrained pattern.  Five MCP servers
sit between the data sources and the triage service:

\begin{itemize}
  \item \texttt{prometheus-mcp:8091}  --- instant and range queries
        against Prometheus.
  \item \texttt{loki-mcp:8092}  --- log queries against Loki.
  \item \texttt{jaeger-mcp:8093}  --- trace lookups against Jaeger,
        plus a \texttt{get\_trace} tool for full-trace retrieval.
  \item \texttt{rca-history-mcp:8094}  --- read access to the
        persistent RCA history (SQLite, on the laptop runner).
  \item \texttt{deploy-events-mcp:8095}  --- deployment-event timeline
        from the application repositories.
\end{itemize}

In the default path, the triage service \emph{itself} (not the LLM)
fans out to these MCP servers during context gathering and assembles
an evidence pack that is rendered into the prompt as a structured
\texttt{\#\# Pre-Gathered Context} section.  The LLM is invoked once
on this pre-gathered material.  Only the bounded-agency retry path
(chapter~\ref{ch:sprint3}, \S\ref{sec:bounded-agency}) gives the LLM
direct access to MCP tools, and only one such tool call is permitted
per retry.

This shape is deliberate: it preserves the LLM's strengths
(structured-output generation, narrative synthesis) while keeping the
operational paths (which queries to run, which evidence to gather)
under deterministic control.  The \textbf{MCP-only invariant} ---
``every datum the LLM sees must arrive via an MCP bridge'' --- is the
project's hallucination firewall in its most general statement.

\section{Related platforms and prior art}

\begin{itemize}
  \item \textbf{Datadog Watchdog} performs anomaly detection across
        metrics and produces narrative incident summaries.  It is a
        proprietary, closed system; the CIRES platform's open-source
        construction allows the team to inspect and tune every layer.
  \item \textbf{New Relic Applied Intelligence} (\emph{NRAI}) provides
        correlation and prioritisation across alerts.  Its
        ``incident intelligence'' graph is a useful precedent for the
        Sprint-4 \emph{incident correlator} (US-5.2) that the CIRES
        platform scopes as future work.
  \item \textbf{Honeycomb BubbleUp} surfaces feature-value
        distributions that distinguish a slow request population from a
        fast one --- an effective bottom-up tool but one that requires
        Honeycomb's specific high-cardinality data model.
  \item \textbf{Komodor} produces narrative incident timelines for
        Kubernetes workloads.  The CIRES platform's exemplar archetypes
        and trace-depth gather (planned via \texttt{get\_trace}) cover
        a similar conceptual space.
\end{itemize}

The closest published precedent for the platform's overall shape is
the body of work around structured-output LLM use in support engineering
(``LLM as RCA copilot'' rather than ``LLM as autonomous responder'').
The novel contributions of this project are (a) the MCP-only invariant
treated as a project-wide rule, (b) the exemplar library of fourteen
canonical RCA archetypes injected as prompt scaffolding, and (c) the
layered prose-shape validator and confidence-clamp mechanism that
together comprise the hallucination firewall.

% =====================================================================
\chapter{Discovery and Requirements (Sprint 0)}
\label{ch:discovery}

\section{Approach}

The project ran a deliberate Sprint~0 between February and the first
days of March 2026 before any infrastructure code was written.  Two
parallel tracks ran for the duration of this sprint:

\begin{enumerate}
  \item \textbf{Monitoring-solution benchmarking.}  A survey of the
        major open-source observability stacks, their feature matrices,
        deployment models, and integration footprint with the JVM /
        Spring Boot stack already in use at CIRES.  SaaS alternatives
        (Datadog, New Relic, Honeycomb) were costed and ruled out for
        this engagement on a mix of budget and data-sovereignty
        grounds.
  \item \textbf{Field research.}  Structured conversations with the
        engineers and operators closest to the production system,
        focused on three questions: what wakes them up, how they
        first discover that something is broken, and what their
        first investigation step is when an alert lands.
\end{enumerate}

A third stream, conducted in parallel, was a guided read-through of
Google's \emph{SRE Book} and \emph{SRE Workbook}, with explicit
translation of the SRE vocabulary into the project's operational
context.

\section{Findings from interviews}

The single most consequential finding from the field research recurred
in nearly every conversation.  When asked how problems were first
discovered, the modal answer was a variant of:

\begin{notebox}[title={Recurring interview finding}]
``We find out things are broken because a client tells us, or because
someone notices after the fact in a log review.''
\end{notebox}

Detection was lagging actual breakage by hours to days.  The team's
reactive posture was structural, not a matter of effort: there was
nothing in the system that drew attention to a problem on its own.
This finding crystallised the problem statement given in
\S\ref{ch:intro} and shaped the entire architecture that followed.

A second finding, less central but still important, was that even
when alerts \emph{were} in place (e.g. for disk usage on certain
servers), the operators routinely had to do non-trivial investigation
work to understand what the alert meant in the current context.  The
project's emphasis on \emph{pre-investigated} alerts --- arriving
with a candidate RCA, evidence, and remediation --- traces back
directly to this second finding.

\section{Solution-space exploration}

Several candidate architectures were brainstormed and pressure-tested
before settling on the eventual two-layer design:

\begin{enumerate}
  \item \textbf{Traditional rules-only alerting.}  Rejected --- it
        already exists conceptually, addresses the
        \emph{detection} half of the problem only, and leaves the
        \emph{investigation} half to a human reading dashboards at the
        right moment.
  \item \textbf{AI-summarisation layer over raw alerts.}  Closer:
        addresses the investigation half partially, by giving the
        operator a narrative summary alongside the alert.  Risky
        because the summary is purely descriptive (``CPU is high'')
        rather than diagnostic (``CPU is high because the cron job is
        running and missing its lock''); it does not actually reduce
        the operator's investigation load.
  \item \textbf{Full AI-triage layer that turns alerts into structured
        RCAs with verdicts and remediation suggestions.}  Chosen.
        Addresses both halves: the observability layer detects, the
        triage layer pre-investigates.  Risks recorded explicitly:
        hallucination, surface-only output, model-availability
        dependence on a self-hosted LLM, latency budget of the LLM
        call vs.\ operator expectations.
\end{enumerate}

\section{Functional requirements}

\begin{enumerate}
  \item \textbf{FR-1} The platform shall ingest metric, log, and trace
        telemetry from the target web stack with no source-code
        modification of the application services.  Instrumentation
        shall be \emph{external}: agents mounted at start time, plug-ins
        configured at the gateway layer.
  \item \textbf{FR-2} The platform shall raise alerts on user-observable
        symptoms (latency, error rate, traffic) and on a set of cause-
        shaped supplementary signals (CPU, memory, disk, target-down,
        log anomalies).
  \item \textbf{FR-3} For every alert that is not suppressed by the
        deduplication or recurrence-gate layers, the platform shall
        produce a structured RCA containing: verdict (escalate /
        dismiss / inconclusive), root-cause prose, suggested actions
        or read-only diagnostic steps, a list of cited evidence, a
        list of correlated alerts in the surrounding window, and a
        confidence value.
  \item \textbf{FR-4} The RCA shall be persisted to a durable history
        store, exposed via a dashboard, and dispatched as an
        escalation email when the verdict so indicates.
  \item \textbf{FR-5} Operator overrides on RCAs (``actually
        dismiss'' / ``actually escalate'') shall be captured and
        used as a routing signal for subsequent fires.
  \item \textbf{FR-6} All data sources consumed by the LLM shall be
        accessed via an MCP bridge (the MCP-only invariant).
\end{enumerate}

\section{Non-functional requirements}

\begin{enumerate}
  \item \textbf{NFR-1 Reproducibility.}  The entire stack shall be
        reproducible from clean infrastructure via a single Ansible
        playbook (Sprint~1) or a Terraform + Ansible pair (Sprint~2).
  \item \textbf{NFR-2 Open-source first.}  No hosted SaaS dependencies
        on the critical path.  An LLM may be a third-party model file,
        but it must be runnable locally on commodity hardware.
  \item \textbf{NFR-3 Latency budget.}  Time from webhook arrival to
        persisted RCA: target $<2$\,min warm, $<3$\,min cold, on the
        laptop runner.  This budget is met (typical 25--68\,s) and
        documented in chapter~\ref{ch:evaluation}.
  \item \textbf{NFR-4 Cost.}  Total monthly run-rate target: under
        \$100/month, exclusive of AWS Free-Tier-eligible resources.
        This budget drives the decision to drop the Sprint~4 GPU
        migration (\S\ref{sec:gpu-dropped}).
  \item \textbf{NFR-5 Auditability.}  Every change to the platform
        shall be reviewable in version control; durable design
        decisions shall be captured in a decisions log; daily
        static and live audits shall produce a written trail.
\end{enumerate}

% =====================================================================
\chapter{System Architecture}
\label{ch:architecture}

\section{Layered architecture overview}

The platform is structured in three layers, illustrated in
figure~\ref{fig:layers}.

\begin{figure}[h]
\centering
\begin{tabular}{|c|}
\hline
\textbf{Layer 3 --- AI triage} \\
FastAPI service, qwen2.5:7b on Ollama, dashboard, SMTP escalation \\
\hline
\textbf{Layer 2 --- MCP bridge} \\
\texttt{prometheus-mcp}, \texttt{loki-mcp}, \texttt{jaeger-mcp},
\texttt{rca-history-mcp}, \texttt{deploy-events-mcp} \\
\hline
\textbf{Layer 1 --- Observability foundation} \\
Prometheus, Loki, Jaeger, Grafana, OTel Collector, node-exporter,
cAdvisor, Drain3 poller \\
\hline
\textbf{Layer 0 --- Telemetry sources} \\
Spring Boot (OTel agent), Kong (OTel plug-in), MySQL, k3s pods,
host metrics \\
\hline
\end{tabular}
\caption{Layered architecture.  Each layer consumes only the layer
immediately below it; the MCP-only invariant is the rule that
formalises this for the AI layer.}
\label{fig:layers}
\end{figure}

The arrangement is deliberately conservative.  Conventional alerts
flow up from layer~0 through Prometheus and Grafana in the standard
way; the AI layer is an additional consumer that does not displace the
existing alert routing.  If the AI layer is taken down, the platform
degrades gracefully to a standard Grafana-alerting setup; the inverse
is not true.

\section{Runtime topology}

After the Sprint~2 migration the platform's runtime topology
distributes across three classes of host:

\begin{enumerate}
  \item \textbf{Monitoring host} (\texttt{observability-rca-monitoring}
        on \texttt{us-east-1}, EC2 \texttt{t3.large}).  Runs the
        observability stack as a seven-container Docker Compose
        deployment: Prometheus, Loki, Jaeger, Grafana, OTel
        Collector, cAdvisor, node-exporter.
  \item \textbf{Application host} (\texttt{observability-rca-k3s} on
        \texttt{us-east-1}, EC2 \texttt{t3.xlarge}).  Runs the
        application workload (Spring Boot, Kong, MySQL, front-end) on
        a single-node k3s cluster.  A node-exporter DaemonSet provides
        host-level metrics from the k3s side.
  \item \textbf{AI host} (the engineer's laptop,
        WSL2 + Docker Desktop, NVIDIA GTX~1060 6\,GB).  Runs the
        triage service, the five MCP servers, and Ollama with
        \texttt{qwen2.5:7b-instruct} loaded.  The laptop is reachable
        over a Tailscale tailnet and addressed by MagicDNS
        (\texttt{adolin-wsl}).
\end{enumerate}

All three hosts share a Tailscale tailnet that also contains the
operator's controller node.  WireGuard handles the transport; the
public IPs of the AWS instances are not exposed for service traffic.

\section{Data flow: alert to RCA}

The end-to-end flow of an alert through the platform proceeds in ten
stages, summarised in figure~\ref{fig:pipeline}.

\begin{figure}[h]
\centering
\small
\begin{tabular}{|r|l|l|}
\hline
\textbf{Stage} & \textbf{What happens} & \textbf{Module} \\
\hline
1  & Webhook arrives from Grafana or Drain3 poller & \texttt{main.py} \\
2  & Two-tier deduplication and suppression        & \texttt{dedup.py} \\
3  & Context gather (Prom, Loki, Jaeger, Drain3, RCA history) & \texttt{context.py} \\
4  & Numeric metric interpretation                 & \texttt{metric\_interpreter.py} \\
5  & Exemplar injection (1 of 14 archetypes)       & \texttt{exemplars/} \\
6  & LLM call (structured output, qwen2.5:7b)      & \texttt{llm\_client.py} \\
6b & Response validation (surface-only, hypothesis-menu, blocklist) & \texttt{response\_validator.py} \\
6c & Bounded-agency retry (one MCP call max)       & \texttt{bounded\_agency.py} \\
6d & Confidence clamp (strip actions if untrustworthy) & \texttt{pipeline.py} \\
7  & Action-template fallback (if RCA actions empty) & \texttt{action\_templates.py} \\
8  & Persist to \texttt{rca\_history}, dispatch email & \texttt{rca\_store.py}, \texttt{notifier.py} \\
\hline
\end{tabular}
\caption{The ten-stage triage pipeline.  Stages 6/6b/6c/6d together
form the hallucination firewall.}
\label{fig:pipeline}
\end{figure}

This pipeline is described stage-by-stage in chapters~\ref{ch:sprint2}
and~\ref{ch:sprint3}; the validator and confidence-clamp mechanisms
introduced in stages~6b and~6d are central to the project's
hallucination-firewall theme and are treated at length in
chapter~\ref{ch:sprint3}.

\section{The MCP-only data-access invariant}

\begin{notebox}[title={MCP-only invariant}]
Every datum that the LLM sees in its prompt must arrive through an
MCP server.  Direct \texttt{httpx} or \texttt{requests} calls to
Prometheus, direct \texttt{aiosqlite} reads of the RCA history
database, and direct in-process \texttt{from app.exemplars import \ldots}
imports are all forbidden in the LLM-adjacent code paths.
\end{notebox}

The invariant serves two purposes.  First, it constrains provenance:
every fact the LLM is shown can be traced back to a specific MCP
tool call, with timestamps and response codes; there is no ambiguity
about ``where did the model see this''.  Second, it removes a class
of hallucination root causes; an in-process Python import can succeed
silently and surface stale or partial data, where an MCP call has an
explicit error envelope and is observable in the bridge layer's
metrics.

Three breaches of the invariant were detected during the
\texttt{0b215ef3} incident on 2026-04-29 and form the scope of the
Epic-4 hallucination-firewall stories (S3-HF-04 / S3-HF-05 /
S3-HF-06).  See chapters~\ref{ch:case-study} and~\ref{ch:sprint3}
for the diagnosis and fix plan.

% =====================================================================
\chapter{Sprint 1 --- Observability Foundation}
\label{ch:sprint1}

\section{Sprint goal and topology}

Sprint~1 ran 2026-03-10 to 2026-03-30.  The goal was to execute on
the first half of the Sprint-0 hypothesis: stand up a full-stack
observability platform for the representative web stack using
open-source tools and deployable by anyone with the playbook.  No
AI layer is present at this stage.

The Sprint-1 topology used three on-premises VMs:
\texttt{monitoring-vm} (.10), \texttt{application-vm} (.11),
\texttt{network-vm} (.12), on subnet 192.168.127.0/24.  All three
were standard Ubuntu LTS hosts with Docker installed; the entirety
of the stack was deployed as Docker Compose services under Ansible
control.

\section{Ansible role layout}

One role per service was established, each idempotent and
parameterised, with a Docker Compose template, environment file
generation, and a health-check task:

\texttt{roles/prometheus}, \texttt{roles/loki}, \texttt{roles/jaeger},
\texttt{roles/grafana}, \texttt{roles/otel\_collector},
\texttt{roles/kong}, \texttt{roles/spring\_boot},
\texttt{roles/cadvisor}, \texttt{roles/node\_exporter},
\texttt{roles/mysql}.

A single playbook (\texttt{playbooks/site.yml}) brought a freshly
provisioned set of VMs to a fully-observable state in approximately
fifteen minutes.  The playbook is part of the project's
reproducibility evidence and is documented in
\texttt{monitoring-docs/reproducibility-guide.html} and
\texttt{build-guide.html}.

\section{Three-pillar telemetry}

\begin{itemize}
  \item \textbf{Metrics} via Prometheus, scraping Spring Boot's
        \texttt{/actuator/prometheus}, Kong's statistics endpoint,
        node-exporter on each VM, and cAdvisor on the
        application-VM.
  \item \textbf{Logs} via Loki, with Promtail tailing Docker container
        logs on each VM and shipping them with a small set of
        well-controlled labels (\texttt{service\_name},
        \texttt{instance}, \texttt{container}).
  \item \textbf{Traces} via the OpenTelemetry Collector, receiving
        spans from the OTel Java agent injected into the Spring Boot
        process at start time and from the OTel plug-in installed at
        the Kong gateway.  Trace identifiers propagated end-to-end so
        a single request appeared as one trace from gateway to
        upstream.
\end{itemize}

\section{Grafana and initial alerting}

Grafana was provisioned via Ansible templates and shipped with four
canonical dashboards (service overview, alert overview, log explorer,
trace explorer) installed under \texttt{/var/lib/grafana/dashboards}.
Initial alerting used Alertmanager, with metric-threshold rules on
CPU, memory, disk, and target-down state, routed through a Gmail SMTP
contact point.  Alertmanager was removed in Sprint~2 (decision~D6) in
favour of Grafana unified alerting with a single \texttt{triage-webhook}
contact point feeding the AI triage service.

\section{Sprint 1 outcome}

\begin{winbox}[title={Sprint 1 deliverables}]
A reproducible, three-pillar observability stack for a representative
web application, deployable from clean VMs in approximately fifteen
minutes via a single Ansible command.  No application source-code
modification required; instrumentation is entirely external (OTel
agent, OTel plug-in, scrape endpoints, Promtail).  This is the
foundation on which all subsequent sprints build.
\end{winbox}

The carry-forward items at sprint close were: the three-VM topology
(superseded in Sprint~2 by AWS EC2 + k3s); Alertmanager (removed in
Sprint~2 by D6); the Gmail SMTP contact point (replaced by the triage
service's own SMTP send); and the on-premises IP address space
(replaced by a Tailscale tailnet).

% =====================================================================
\chapter{Sprint 2 --- AWS Migration and AI RCA MVP}
\label{ch:sprint2}

\section{Sprint goal}

Sprint~2 ran 2026-03-31 to 2026-04-23, with a follow-on Epic~5 running
to 2026-04-29.  The goal was to move the platform off on-premises VMs
onto AWS (Terraform-provisioned EC2), migrate the application workload
to k3s, and add an AI layer that produces structured root-cause analyses
from incoming alerts.  The on-premises foundation from Sprint~1 remains
useful as a reproducibility target; the production-shaped deployment
lives on AWS.

The sprint deadline moved twice (April 9 $\rightarrow$ April 23),
captured in the strikethrough trail in the project's progress tracker.
Final landing: 2026-04-23.

\section{Epic 1 --- AWS Infrastructure}

A Terraform configuration in the \texttt{provisioning-monitoring-infra}
repository manages the entire AWS estate: a VPC with public and
private subnets, two EC2 instances (\texttt{t3.large} for monitoring,
\texttt{t3.xlarge} for k3s), an RDS \texttt{db.t3.micro} MySQL
instance, Elastic IPs on the durable hosts, security groups
parameterised by Tailscale CIDR plus a public-UI CIDR for Grafana /
Prometheus / Loki / Jaeger UI access, IAM roles per instance, an S3
bucket for Loki cold storage and a second for Terraform state.  The
Ubuntu AMI was pinned to a specific image identifier to remove silent
drift between runs.

\section{Epic 2 --- AI RCA System MVP}

This epic introduces the largest new surface in the project.  It
comprises:

\begin{itemize}
  \item \textbf{Five MCP servers}
        (\texttt{monitoring-mcp-servers/}).  Each is a small Python
        process that exposes three to four tools over a JSON-RPC
        transport with a Prometheus-style \texttt{/metrics} endpoint
        for observability of the bridge itself.
  \item \textbf{FastAPI triage service}
        (\texttt{monitoring-triage-service/}).  Webhook receivers for
        Grafana and Drain3, an asynchronous ten-stage pipeline,
        structured-output LLM calls via Ollama, deduplication and
        suppression, a dashboard, and SMTP escalation.
  \item \textbf{Cross-pillar correlation.}  The \texttt{context.py}
        module fans out to three MCPs in parallel before the LLM call,
        building one pre-gathered evidence pack the LLM consumes.  This
        keeps the model out of the loop on which queries to run; the
        platform decides.
  \item \textbf{Drain3 log-template clustering.}  Server-side over
        Loki, surfacing novel templates as a fourth signal alongside
        metrics, logs, and traces.
  \item \textbf{Local LLM.}  \texttt{qwen2.5:7b-instruct} via Ollama,
        initially on a \texttt{g4dn.xlarge} EC2 (decommissioned), then
        moved to the laptop's GTX 1060 on 2026-04-23 (still the
        production runner today).
\end{itemize}

\section{Epic 3 --- Kubernetes migration}

The application workload was moved from the application-VM Docker
Compose into a single-node k3s cluster on
\texttt{observability-rca-k3s}.  Helm charts were authored for the
Spring Boot back-end, Kong, and the front-end.  The
\texttt{xanmanning.k3s} Ansible Galaxy role provisions k3s as a
\texttt{systemd} unit; a node-exporter DaemonSet provides host-level
metrics from the k3s side.

The monitoring stack stayed on Docker Compose.  Decision~D11
records the rationale: the monitoring VM is operationally stable,
deployable from Ansible, and gains nothing structural from migrating
onto Kubernetes; the cost of the migration (Helm charts for nine
services, persistent-volume claim management, Kubernetes-specific
operational learning curve) exceeded the benefit.

\section{Epic 4 --- AI Triage Hardening}

This epic raises the AI triage layer from ``MVP that responds to
webhooks'' to ``service operators can trust'': a two-tier
deduplication layer (in-memory fingerprint plus a pre-LLM
history-based suppression --- decision~D4), structured-output schema
enforcement at the LLM call layer (D12), an error envelope on every
MCP response, pipeline timeouts, the operator dashboard, the email
notifier with a confidence-bordered card, and a first pass at
action-template fallback.

\section{Epic 5 --- UEBA-flavoured stories (post-sprint)}

Five UEBA-inspired stories ran past sprint close as a coherent
in-flight workstream and were ultimately folded into Sprint~3 as the
``RCA Quality v2 close-out'' epic.  The net effort was approximately
forty-nine story points shipped after sprint close.  These stories
are summarised in table~\ref{tab:epic5}.

\begin{table}[h]
\centering
\small
\begin{tabular}{p{0.10\textwidth} p{0.55\textwidth} p{0.25\textwidth}}
\toprule
\textbf{Story} & \textbf{What it added} & \textbf{Outcome} \\
\midrule
US-5.1 Phase A & Per-service Drain3 --- dictionary of \texttt{service
$\to$ TemplateMiner} with file-backed state per service. & Shipped
2026-04-28 \\
US-5.1 Phase B & Entity baselines --- 7-day rolling per-service-metric
quantiles + $\sigma$-claim helper. & Shipped 2026-04-28; verification
closed by S3-HF-04 \\
US-5.3 & Closed-loop feedback override/confirm + lazy
\texttt{triage\_precision} / \texttt{triage\_recall} metrics. & Shipped
2026-04-28 \\
US-5.4 & Adaptive thresholds on three flappiest rules
(same-hour-7-days-ago + $3\sigma$ baseline). & Shipped 2026-04-27 \\
US-5.8 & Recurrence gates (pre-LLM + post-LLM, Medium-CPU/Mem opt-in).
& Shipped 2026-04-28 \\
\bottomrule
\end{tabular}
\caption{Epic 5 stories carried past Sprint 2 close.}
\label{tab:epic5}
\end{table}

\section{Sprint 2 outcome}

\begin{winbox}[title={Sprint 2 deliverables}]
The transition from ``we have observability'' (Sprint~1) to ``the
observability platform produces structured, validated RCAs end-to-end
with operator-feedback closure'' (Sprint~2 + Epic~5).  The pipeline
that runs in production today is the Sprint-2 architecture with
Sprint-3's quality and safety mechanisms layered on top.  The test
suite reached 178/178 at sprint close on the triage service.
\end{winbox}

% =====================================================================
\chapter{Sprint 3 --- Quality, Feedback, and Hallucination Firewall}
\label{ch:sprint3}

\section{Sprint shape}

Sprint~3 runs 2026-04-23 to 2026-05-19 (extended from the original
05-08 close).  Four epics:

\begin{enumerate}
  \item \textbf{Epic 1 --- RCA Quality v2 close-out (49 SP).}
        Backfill of the Epic-5 work that shipped between 2026-04-23
        and 2026-04-29 under Sprint-3's epic record.  Twelve of
        thirteen stories Done; one cleanup story remains.
  \item \textbf{Epic 2 --- Operator feedback loop (7 SP).}  Schema
        extension on the existing US-5.3 feedback table; new
        endpoints; four new MCP tools on \texttt{rca\_history\_mcp};
        an in-email rating form.  Not started at the time of writing.
  \item \textbf{Epic 3 --- Drain3 baseline lifecycle (6 SP).}  Wire
        the dormant S3 plumbing in \texttt{drain\_analyzer.py};
        APScheduler weekly snapshot; boot-time restore; three
        operator endpoints; IAM update.  Not started at the time of
        writing.
  \item \textbf{Epic 4 --- Hallucination firewall + trace depth
        (15 SP, closed 2026-05-19).}  Added on 2026-04-29 after the
        \texttt{0b215ef3} incident.  Two themes: (a) close the three
        MCP-only-invariant bypasses on the LLM-gathering path; (b)
        wire the \texttt{get\_trace(trace\_id)} MCP tool that
        existed but was never invoked.  Five SP shipped same-day as
        the incident; the remaining ten SP (S3-HF-04 through
        S3-HF-09) shipped in one focused session at sprint close-out
        on 2026-05-19.
\end{enumerate}

The sprint contains a deliberate twenty-day investigation window
(2026-04-30 to 2026-05-19) in the middle of its calendar; the work
done during that window is described in \S\ref{sec:investigation}.

\section{The 4-axis RCA quality rubric}

To make ``good RCA'' tractable as an engineering target, the project
defines a four-axis rubric used to score every produced RCA, both in
chaos-run automation and in the daily live audit:

\begin{enumerate}
  \item \textbf{Cause-named} --- the RCA names a root cause, not just
        a restatement of the alert (e.g.\ ``connection pool exhaustion
        on \texttt{spring-boot}'' rather than ``the alert indicates
        high latency'').
  \item \textbf{Evidence-cited} --- the cause is supported by specific
        metric values, log lines, or trace attributes drawn from the
        pre-gathered context.
  \item \textbf{Actionable} --- the RCA carries either operationally
        valid suggested actions or, in the suppressed case, valid
        read-only diagnostic steps.
  \item \textbf{Prose-shape} --- the RCA leads with the cause rather
        than with a PromQL expression, observation, or
        ``based on context'' opener (the philosophy stated by
        decision~D19).
\end{enumerate}

Each axis is scored on a 0--1 scale; the unweighted mean is the
overall RCA quality.  The rubric is implemented in
\texttt{monitoring-project/scripts/chaos/lib/rca\_scorer.py} and used
both by the chaos-injection harness and the daily live audit.

\section{The exemplar library (D17)}

\begin{notebox}[title={Decision D17 --- Exemplar injection}]
The LLM is shown one of fourteen canonical RCA archetypes as a
structural target before it sees the evidence pack.  Each archetype
contains an alertname pattern, a service-and-deployment-type filter,
a signal pillar (metric / log / trace / drain3 / mixed), a severity
class, and a short example RCA in the project's preferred prose
shape.
\end{notebox}

The fourteen archetypes are: \texttt{oom-loop},
\texttt{upstream-latency-attribution},
\texttt{service-self-p95-saturation},
\texttt{critical-resource-saturation},
\texttt{otel-collector-degradation},
\texttt{synthetic-blip-dismiss},
\texttt{cascade-incident-disk-loki},
\texttt{drain3-novelty-post-deploy},
\texttt{bounded-agency-resolves-inconclusive},
\texttt{adaptive-threshold-noop},
\texttt{closed-loop-feedback-override},
\texttt{crashloop-bad-config},
\texttt{network-firewall-attribution},
\texttt{tls-cert-expiry-pre-failure}.  A default
\texttt{generic-sre-shape} fallback is used when none match.

Selection is via
\texttt{exemplars.find\_for\_alert(alertname, service, deployment\_type,
signal, severity)} which scores all fourteen on a weighted match
(alertname regex 0.10, services 0.40, deployment-type 0.30, signal
0.20, severity 0.10) and returns the highest.  The matched exemplar
is rendered as a \texttt{\#\# Reference exemplar} section
\emph{before} the \texttt{\#\# Pre-Gathered Context} in the prompt ---
the placement is load-bearing and was confirmed during the
investigation window (\S\ref{sec:investigation}) to outperform
post-context placement on the 4-axis rubric.

\section{Adaptive thresholds (D18)}

Static numeric thresholds (\texttt{p95 > 1000\,ms},
\texttt{cpu > 80\,\%}) generate floods of false positives during
off-hours and miss real anomalies during peak traffic.  Real traffic
isn't stationary; rules shouldn't pretend it is.  Three Grafana rules
(\texttt{HighP95Latency}, \texttt{HighKongP95Latency},
\texttt{MediumCpuUsage}) were converted to a same-hour-7-days-ago
baseline plus a $3\sigma$ envelope.  The thresholds adjust on a
weekly rolling window and the alert fires only when the current
value lies outside the historical envelope for the matching hour.

\section{RCA prose philosophy (D19)}

Decision~D19 captures the position that \emph{telemetry is the means,
not the end}: an RCA paragraph that leads with a PromQL expression
or an observed metric value is reading like the alert query, not like
an investigation.  Three reinforcing surfaces were updated to enforce
a cause-first prose shape:

\begin{itemize}
  \item \texttt{SYSTEM\_PROMPT} rule~A and a new rule~J explicitly
        instruct the model to lead with the cause; the previous version
        of rule~A had instructed the model to lead with the PromQL
        expression and observed value, an unintended training surface
        for surface-only output.
  \item All eleven exemplar \texttt{rca} paragraphs were rewritten to
        match the cause-first shape.
  \item A new validator rule scans the produced RCA against a set of
        \texttt{\_SURFACE\_ONLY\_LEDE\_PATTERNS} regexes.  Matches
        trigger a retry under the bounded-agency path.
\end{itemize}

\section{Closed-loop feedback (D20)}

Operator overrides on RCAs were previously silent: the LLM produced a
verdict; the operator either agreed or didn't; the system never
learned.  Without that signal, every prompt, exemplar or threshold
change was unmeasurable.

Decision~D20 introduces a closed-loop layer.  A \texttt{feedback}
table records the operator's actual verdict ($\pm$ a free-text
\texttt{actual\_root\_cause}).  Two endpoints expose the surface
(\texttt{/feedback/override} and \texttt{/feedback/confirm}).  Two
Prometheus metrics, \texttt{triage\_precision} and
\texttt{triage\_recall}, are computed lazily from the feedback
table.  And the suppression layer takes operator overrides into
account: an alert with a recent ``actually dismiss'' override is
suppressed pre-LLM the next time it fires.

\section{Recurrence gates (D21)}

An alert that fires every ninety minutes is too slow for fingerprint
deduplication (10-minute window) but too noisy to page someone every
time.  Decision~D21 introduces a recurrence gate at two points in the
pipeline: a pre-LLM gate that consults the RCA history for prior
fires of the same alert key, and a post-LLM gate that consults the
operator override layer from D20.  Critical-severity alerts bypass
the gate.  Per-rule opt-in is via a Grafana rule annotation
\texttt{recurrence\_gate}; \texttt{MediumCpuUsage} and
\texttt{MediumMemoryUsage} are currently opted in.

\section{Per-service Drain3 template trees (D23)}

The original Drain3 integration used a single global
\texttt{TemplateMiner}.  Spring Boot's ``novel'' templates were
judged against the global pool, not against Spring Boot's own
history.  A line that's commonplace for \texttt{spring-boot} but
unusual for \texttt{kong} looked the same as a line that's unusual
for both.  Decision~D23 introduces a per-service split: a dictionary
\texttt{dict[service $\to$ TemplateMiner]} with \texttt{FilePersistence}
per service, and a migration path that converts the pre-split
\texttt{drain3\_state.bin} into a service-keyed family.

\section{The hallucination firewall epic (Epic 4)}

Epic~4 was added on 2026-04-29 after the \texttt{0b215ef3} incident
(chapter~\ref{ch:case-study}).  The case-study chapter covers the
incident itself; this section summarises the resulting stories:

\begin{itemize}
  \item \textbf{S3-HF-01 (Tier 0, 1 SP, shipped 2026-04-29).}
        Confidence clamp strips \texttt{suggested\_actions} when the
        clamp fires and emits a read-only
        \texttt{diagnostic\_steps} field instead.  Email and dashboard
        render the diagnostic steps as a separate card with an
        explicit ``Withheld --- confidence clamped to 0.40'' notice.
  \item \textbf{S3-HF-02 (1 SP, shipped 2026-04-29).}  Five defensive
        tests pinning the action-template lookup against future
        regex-widening regressions.
  \item \textbf{S3-HF-03 (Tier 2, 2 SP, shipped 2026-04-29).}
        Hypothesis-menu validator: seven regex patterns targeting
        ``possibly $X$ or $Y$'' / ``may be due to (slow query | pool |
        GC)'' / ``either $A$ or $B$'' prose with negative lookahead
        for idiomatic ``either way'' / ``either case''.  Companion
        cause-evidence rule that tokenises both RCA and evidence with
        \texttt{[a-z]\{4,\}} (so
        \texttt{hikaricp\_connections\_active} splits into
        \texttt{[hikaricp, connections, active]} for cross-matching)
        and requires non-stop-word overlap.
  \item \textbf{S3-HF-04 (1 SP, shipped 2026-05-19).}  Route
        \texttt{entity\_baselines.py:144} through
        \texttt{prometheus-mcp} instead of direct \texttt{httpx} to
        Prometheus \texttt{/api/v1/query}.  Closes one of three
        MCP-only-invariant bypasses.
  \item \textbf{S3-HF-05 (1 SP, shipped 2026-05-19).}  Route the
        \texttt{bounded\_agency.py} RCA-history paths through
        \texttt{rca-history-mcp} instead of direct \texttt{aiosqlite}.
  \item \textbf{S3-HF-06 (2 SP, shipped 2026-05-19).}  Extend
        \texttt{rca\_history\_mcp} with quality-rated tools:
        \texttt{get\_similar\_decisions(min\_quality)} and
        \texttt{get\_low\_rated\_examples\_for\_alert}.  Eight new
        MCP-level tests cover the rating-filter semantics.
  \item \textbf{S3-HF-07 (Tier 1, 3 SP, shipped 2026-05-19).}  Deep
        trace gather via \texttt{get\_trace(trace\_id)} MCP.  The
        pipeline previously gathered Jaeger data at the
        service-boundary level only and could identify ``upstream is
        slow'' but not name what was slow.  With trace depth, the LLM
        now sees per-span attributes (\texttt{db.statement},
        \texttt{http.target}, error tags) for the failing trace.
        Pre-step (executed during the investigation window): sampled
        12 live Spring Boot traces to confirm JDBC parameterisation;
        \texttt{db.statement} spans use bind parameters
        (\texttt{?,\,?,\,?}), so PII redaction was not a blocker.
  \item \textbf{S3-HF-08 (2 SP, shipped 2026-05-19).}
        MCP-integrity continuous-integration lint forbidding
        \texttt{httpx}/\texttt{requests}/direct-DB calls outside
        \texttt{\_mcp\_call} helpers and the MCP servers themselves.
        The lint runs on every push and was the work that confirmed
        S3-HF-04 and S3-HF-05 were the only two bypasses on the
        LLM-gather path.
  \item \textbf{S3-HF-09 (Tier 4, 2 SP, shipped 2026-05-19).}  Chaos
        scorer 5th axis, labelled corpus seed (the \texttt{0b215ef3}
        incident itself), and a CI regression gate.  Landed before
        Sprint~4's Tier~3 iterative agentic gather, which is now
        unblocked by the Tier~4 measurement scoreboard being in
        place.
\end{itemize}

Epic~4 closed at \textbf{15/15 SP on 2026-05-19} in a single focused
session.  All three MCP-only-invariant bypasses on the
LLM-gather path are now closed, the trace-depth path is wired, the CI
lint forbids future regressions, and the chaos corpus regression gate
is live (suite \texttt{252/252}).

\section{The 20-day investigation window (2026-04-30 to 2026-05-19)}
\label{sec:investigation}

After Epic~4 was added on 2026-04-29 and the Tier-0/2 stories shipped
the same day, the team deliberately slowed down before the larger MCP
refactors (S3-HF-04 through S3-HF-09).  Approximately three weeks
went into prototyping output-quality improvements in scratch branches
and refreshing documentation, with no commits to \texttt{main} by
design --- the goal was to know which ideas were worth the refactor
cost before paying it.  Nothing structural shipped during the window;
the production triage container stayed on commit \texttt{d6d10ef} from
2026-04-29, and the seventeen daily static audits all reported the
test suite green at 226/226 with the same set of seven minor
doc-drift items.  Eight candidate output-quality improvements were
A/B-tested against a twenty-eight-decision sample drawn from the
live-audit log and chaos runs 3 and 4; the outcomes are summarised in
table~\ref{tab:investigation}.

\begin{table}[h]
\centering
\footnotesize
\begin{tabular}{p{0.32\textwidth} p{0.40\textwidth} p{0.18\textwidth}}
\toprule
\textbf{Idea} & \textbf{Result} & \textbf{Decision} \\
\midrule
Exemplar / context section ordering --- exemplar after context. &
Cause-named score dropped $0.71 \to 0.62$ on validator-flagged subset; the
LLM dropped the structural cue once the evidence was in scope. &
Keep pre-context order \\
\midrule
Evidence-pillar weighted re-ranker. & No measurable effect; RCAs
already lead with the firing pillar in 24/28 cases. & Shelved \\
\midrule
BM25 retrieval prototype. & Degenerate over fourteen archetypes;
top-1 hit always matched
\texttt{exemplars.find\_for\_alert}.  Sprint-4 candidate \#13 stays
as scoped. & Confirms retriever choice is not the gate \\
\midrule
Hypothesis-menu validator threshold sweep. & 0/120 false positives at
strict mode; warn-only adds no catches. & Default strict stays \\
\midrule
Confidence-clamp threshold (0.30 / 0.40 / 0.50). & Below 0.40 dampens
prose; above 0.40 actions leak through.  & 0.40 retained \\
\midrule
Trace-depth feasibility scan on 12 live Spring Boot spans. &
\texttt{db.statement} uses bind parameters not literal values. &
S3-HF-07 has no PII blocker \\
\midrule
RCA prose A/B (cause-first vs.\ observation-first). & Observation-first
regresses on cause-named when evidence is thin. & Cause-first retained
(reinforces D19) \\
\midrule
Pipeline-duration per-stage instrumentation. & Drafted in scratch
branch. & Did not merge; depends on S3-HF-05 \\
\bottomrule
\end{tabular}
\caption{Investigation-window outcomes (2026-04-30 to 2026-05-19).
Each candidate was scored on the 4-axis RCA quality rubric over a
twenty-eight-decision sample.}
\label{tab:investigation}
\end{table}

\section{Documentation refresh}

A substantial documentation refresh landed on 2026-05-19 as the
window closed:

\begin{itemize}
  \item A new canonical \texttt{sprint-history.html} consolidating
        Sprints 0--4 onto one page (there had been no per-sprint
        timeline before).
  \item A new \texttt{solution-brief.html} as a report-facing
        one-pager, built around the \texttt{0b215ef3} case study, the
        validator details, the fourteen-archetype table, and a
        glossary.
  \item Metric-name doc-vs-code mismatches reconciled in
        \texttt{triage-service.html} \S7.
  \item Consistency passes across \texttt{architecture.html} and
        \texttt{system-explained.html}.
  \item Tested-and-rejected experiments folded back as footnotes to
        D17 / D19 / D21 in \texttt{decisions-log.html} so the
        prototyping outcomes are durably recorded next to the
        decisions they tested.
\end{itemize}

\section{The bounded-agency retry}
\label{sec:bounded-agency}

The bounded-agency retry is the platform's mechanism for handling
genuinely data-starved RCAs: if the validator flags the first pass,
the LLM is given exactly one additional MCP call from a six-tool
whitelist (Prometheus instant query, Loki range, Jaeger get-traces,
RCA-history lookup, deploy-events, and \texttt{get\_trace}).  The
whitelist plus the Pydantic schema plus a retry feedback string
constrain the second call; the model is re-prompted, the output is
revalidated, and the first-pass output is kept if the retry is still
bad.  This is decision~D15, ``one MCP call, not a chain'' --- a
deliberately constrained version of agentic gather chosen for
predictability and bounded latency, with the unbounded
\emph{iterative agentic gather} scoped to Sprint~4 once the Tier-4
measurement scoreboard is in place.

% =====================================================================
\chapter{Case Study: the 0b215ef3 Incident}
\label{ch:case-study}

\section{The alert}

On 2026-04-29, a \texttt{HighKongP95Latency} alert fired for the
\texttt{spring-boot} upstream.  The Kong gateway's p95 latency had
crossed its adaptive threshold for the second time that hour.  The
webhook arrived at the triage service on the laptop at
\texttt{adolin-wsl} and produced the RCA with identifier
\texttt{0b215ef3}.

\section{The bad RCA}

The produced RCA carried:

\begin{itemize}
  \item \textbf{Verdict.} \texttt{escalate}.
  \item \textbf{Confidence.} \texttt{0.55} (moderate).
  \item \textbf{Root-cause prose.}  A hedged ``hypothesis menu'':
        ``This could be a slow query in the
        \texttt{products} service, a connection-pool exhaustion event,
        or a garbage-collection pause; the latency profile is
        consistent with any of the three.''
  \item \textbf{Suggested actions.}  Three templated \texttt{kubectl}
        commands operating on a \texttt{products} deployment that does
        not exist (the system runs Spring Boot, not a microservice
        called \texttt{products}).
  \item \textbf{Evidence.}  Two Prometheus queries against
        \texttt{kong\_http\_requests\_total} with their text rendered
        but no values cited; one trace identifier without
        \texttt{db.statement} or per-span attributes.
\end{itemize}

\begin{warnbox}[title={Why the operator caught it}]
The hypothesis menu was the visible symptom.  The deeper failures
were structural: the LLM had been given the latitude to invent a
\texttt{products} service that does not exist; the confidence value
was not low enough to trip the existing clamp at 0.40; and the
evidence pack was too thin for the LLM to commit to a single root
cause.
\end{warnbox}

\section{Forensic analysis: three MCP-invariant violations}

The post-incident analysis identified three distinct code paths that
violated the MCP-only invariant and contributed to the bad output:

\begin{enumerate}
  \item \texttt{monitoring-triage-service/app/entity\_baselines.py:144}
        was using a direct \texttt{httpx} call to Prometheus's
        \texttt{/api/v1/query} endpoint, bypassing
        \texttt{prometheus-mcp}.  The entity-baseline values for the
        Kong p95 metric flowed into the context pack with no MCP
        provenance and (it transpired) a slightly stale window.
  \item \texttt{monitoring-triage-service/app/bounded\_agency.py}
        was using direct \texttt{aiosqlite} reads of the
        \texttt{rca\_history} table for retry-tool dispatch,
        bypassing \texttt{rca-history-mcp}.  This meant the retry path
        could not benefit from the
        \texttt{get\_similar\_decisions(min\_quality)} filter that
        would have surfaced operator-confirmed similar cases.
  \item \texttt{monitoring-triage-service/app/llm\_client.py}
        and the exemplar consumers used in-process
        \texttt{from app.exemplars import \ldots} imports rather than an
        MCP tool.  In itself this is benign, but it left the exemplar
        layer outside the project's provenance regime and would have
        broken when the same triage service moved to a multi-process
        deployment.
\end{enumerate}

A fourth, related, finding was that the existing
\texttt{get\_trace(trace\_id)} MCP tool --- written months earlier ---
had never been called from the pipeline.  The context gather was
operating at the service-boundary level: it knew ``upstream is slow''
but could not name what was slow because it had not pulled the spans
of the offending trace.

\section{The tiered response}

A planning agent produced a five-tier response plan, captured in the
session handoff and translated into the Epic~4 stories listed in
\S\ref{ch:sprint3}.  In summary:

\begin{itemize}
  \item \textbf{Tier 0} --- strip suggested actions when the confidence
        clamp fires; emit read-only diagnostic steps instead.
        Shipped same day as S3-HF-01 (2026-04-29).
  \item \textbf{Tier 1} --- deep trace gather via
        \texttt{get\_trace(trace\_id)}.  Shipped as S3-HF-07 on
        2026-05-19.
  \item \textbf{Tier 2} --- hypothesis-menu validator and
        cause-evidence rule.  Shipped same day as S3-HF-03
        (2026-04-29).
  \item \textbf{Tier 3} --- iterative agentic gather with
        hypothesis-tree state.  Deferred to Sprint~4; the Tier~4
        scoreboard dependency is now satisfied (S3-HF-09 shipped),
        so the Sprint-4 gate is operational rather than structural.
  \item \textbf{Tier 4} --- chaos scorer 5th axis, corpus seed,
        regression gate.  Shipped as S3-HF-09 on 2026-05-19.
\end{itemize}

The same-day work shipped as commits \texttt{b0f1d0c} and
\texttt{d6d10ef} on the triage service, plus accompanying commits on
the monitoring-docs and Jira-CSV repositories.  The test suite went
from 185 mid-day to 226/226 at end-of-day --- forty-one new tests
across the three shipped stories, with zero regressions.

\section{Lessons captured}

The incident is the canonical case study in the project for two
reasons.  First, it exposed three structural violations of a stated
invariant; the platform's documentation had claimed an invariant that
the code did not enforce.  This is the worst class of documentation
drift, and the discovery prompted both the Epic-4 stories and the
S3-HF-08 CI lint that will prevent future violations.  Second, the
incident validated the confidence-clamp design choice: the clamp did
fire, but at 0.55 the configured ceiling of 0.40 was high enough that
the bad actions leaked through.  The clamp behaviour was strengthened
(strip actions, emit diagnostic steps) rather than re-thresholded;
the investigation-window threshold sweep
(table~\ref{tab:investigation}) confirmed 0.40 was the right ceiling.

% =====================================================================
\chapter{Evaluation}
\label{ch:evaluation}

\section{Test-coverage progression}

The triage-service test suite size over the project's lifetime is a
useful coarse measure of the engineering load behind the visible
features.  Table~\ref{tab:tests} summarises the milestones.

\begin{table}[h]
\centering
\small
\begin{tabular}{lrl}
\toprule
\textbf{Date} & \textbf{Suite size} & \textbf{Milestone} \\
\midrule
2026-04-27 evening & 89  & Exemplar library (D17) tests in place \\
2026-04-28 mid-day & 95  & RCA-prose validator (D19) tests added \\
2026-04-28 EOD    & 178 & Full RCA quality chain (F-1 to F-5) live \\
2026-04-29 EOD    & 226 & S3-HF-01/02/03 shipped \\
2026-05-19        & 226 & Held green across 17 daily audits \\
\bottomrule
\end{tabular}
\caption{Triage-service test-suite progression.}
\label{tab:tests}
\end{table}

The point at which the suite size more than doubled in a single
24-hour window (95 $\to$ 178) marks the F-1 to F-4 RCA-quality work
described in chapter~\ref{ch:sprint2}.  The +48 over the
2026-04-29 day (178 $\to$ 226) is the Epic-4 same-day stories.

\section{RCA quality on the 28-decision sample}

The investigation-window evaluation produced a baseline RCA quality
score on a 28-decision sample drawn from the live-audit log and
chaos runs 3 and 4.  On the 4-axis rubric:

\begin{table}[h]
\centering
\small
\begin{tabular}{lr}
\toprule
\textbf{Axis} & \textbf{Mean score} \\
\midrule
Cause-named       & 0.78 \\
Evidence-cited    & 0.84 \\
Actionable        & 0.71 \\
Prose-shape       & 0.89 \\
\midrule
Overall (unweighted) & 0.81 \\
\bottomrule
\end{tabular}
\caption{Baseline RCA quality on the 28-decision sample
(higher is better).  Actionable is the lowest score because of the
deliberate strip-actions-on-clamp behaviour (S3-HF-01), which lowers
the actionability score but raises operational trust.}
\label{tab:rca-quality}
\end{table}

\section{Critical RCA output audit --- 12 representative cases (2026-05-21)}
\label{sec:critical-rca-audit}

The 28-decision baseline in the previous section uses a permissive
four-axis rubric inherited from the chaos harness.  A complementary
critical audit ran 2026-05-21 against the persisted
\texttt{/decisions} corpus (271 rows over the 28-day operating
window), sampling 12~RCAs across the six dominant alert families
(\texttt{TargetDown}, \texttt{HighKongP95Latency},
\texttt{MediumCpuUsage}, \texttt{HighP95Latency},
\texttt{Drain3AnomalyDetected}, \texttt{PodHighMemoryUsage}) with
deliberate inclusion of both verdict outcomes per family.  The
purpose was to evaluate the platform against the stricter standards
the P0--P11 Sprint-4 plan was scoped to enforce, not the looser
chaos-harness baseline.

\subsection{Audit rubric}

Eight pass/fail criteria, grounded in the project's documented
output-quality standards
(\texttt{feedback\_rca\_prose\_quality} + the P0--P11
findings):

\begin{enumerate}
  \item \textbf{C1 Cause-first lede} --- first sentence of
        \texttt{rca\_report} names a specific root cause or
        mechanism, not the alert.
  \item \textbf{C2 Verdict appropriateness} ---
        \texttt{escalate}/\texttt{dismiss} defensible from the
        available evidence.
  \item \textbf{C3 Evidence concreteness} --- \texttt{evidence}
        list contains specific values, not template prose or
        placeholders.
  \item \textbf{C4 No-noise prose (P1)} --- no PromQL function
        calls, raw floats $>$4 significant figures, Z-score
        equations, or \texttt{\# TYPE}/\texttt{\# HELP}
        Prometheus-output markers in \texttt{rca\_report}.
  \item \textbf{C5 Suggested actions correct (P5)} --- actions
        are concrete + deployment-type-correct; empty list is
        acceptable only for \texttt{dismiss} verdicts.
  \item \textbf{C6 Diagnostic-steps non-templated (P3)} ---
        same-alert-type RCAs from different fires have $\geq$3
        distinct diagnostic-step strings each.
  \item \textbf{C7 \texttt{rca\_quality} tag accurate (P11)} ---
        the \texttt{actionable} tag requires (named cause)
        $\land$ ($\geq$1 evidence item) $\land$ ($\geq$1 action
        OR explicit dismiss reasoning).
  \item \textbf{C8 No truncated-metric hallucination (P4)} ---
        \texttt{rca\_report} does not cite truncation artefacts
        (e.g. metric names ending in
        \texttt{\_daem (truncated)\ldots]}) as real metric
        names.
\end{enumerate}

\subsection{Per-criterion pass rate}

\begin{table}[h]
\centering
\small
\begin{tabular}{llrr}
\toprule
\textbf{ID} & \textbf{Criterion} & \textbf{Pass} & \textbf{Fail/borderline} \\
\midrule
C1 & Cause-first lede           &  6/12 (50\%) & 6/12 \\
C2 & Verdict appropriateness    &  7/12 (58\%) & 5/12 \\
C3 & Evidence concreteness      &  5/12 (42\%) & 7/12 \\
C4 & No-noise prose (P1)        &  4/12 (33\%) & 8/12 \\
C5 & Actions correct (P5)       &  5/12 (42\%) & 7/12 \\
C6 & Diagnostic-steps non-template (P3) & 2/12 (17\%) & 10/12 \\
C7 & \texttt{rca\_quality} tag accurate (P11) & 4/12 (33\%) & 8/12 \\
C8 & No truncated-metric (P4)   &  9/12 (75\%) & 3/12 \\
\midrule
\textbf{Cases rated GOOD overall}     & & 3/12 (25\%) & \\
\textbf{Cases rated OK overall}       & & 2/12 (17\%) & \\
\textbf{Cases rated BAD overall}      & & 7/12 (58\%) & \\
\bottomrule
\end{tabular}
\caption{Critical RCA audit per-criterion pass rates on 12 sampled
cases.  C4 (no-noise) and C6 (non-templated diagnostic steps) are
the two worst dimensions; C8 (no truncated-metric hallucination) is
the strongest only because the specific artefact is rare in the
corpus, not because the safeguard exists yet (P4 is unshipped).}
\label{tab:critical-rca}
\end{table}

\subsection{Failure-mode taxonomy with specific cases}

Six dominant failure modes recur across the sample, in descending
order of severity.

\paragraph{F1.~Exemplar placeholders leak through verbatim (P6).}
Cases \texttt{1a95f6a6} (\texttt{HighKongP95Latency},
\texttt{confidence=0.92}) and \texttt{0c76258b}
(\texttt{Drain3AnomalyDetected}, \texttt{confidence=0.82}) emitted
RCA prose containing literal \texttt{\{service\}}, \texttt{<X>},
\texttt{<verbatim>}, \texttt{Y-Z}, \texttt{HH:MM} placeholders
that the LLM failed to substitute with concrete values.  Both
carry high LLM confidence, which makes the failure dangerous: the
output looks authoritative but the substitution layer never ran.
This is the textbook case for Sprint~4 P6
(\texttt{S4-PR-06} exemplar prose rewrites + overlap-detection
validator).

\paragraph{F2.~PromQL / numerical noise in cause-prose (P1).}
Cases \texttt{a9393ef6} (full \texttt{up\{instance=...\} = 0}
PromQL in the lede), \texttt{0b215ef3} (\texttt{histogram\_quantile(0.95)
= 8203.846153846154 ms}), \texttt{79119df0} (the full
\texttt{100 - (avg by(instance) (rate(node\_cpu\_seconds\_total\{mode="idle"\}[1m])) * 100)}
expression in the lede), and \texttt{d763a9c9} (\texttt{\# TYPE
jvm\_threads\_daem (truncated)\ldots]} quoted verbatim).  Four of
twelve sample RCAs leak raw telemetry into operator-facing prose.
Sprint~4 P1 (\texttt{S4-PR-01}) scrubber is scoped for this.

\paragraph{F3.~\texttt{rca\_quality=actionable} over-applied (P11).}
Seven of twelve sampled RCAs carry \texttt{actionable} when they
shouldn't.  \texttt{0b215ef3} has actions but they are
wrong-archetype (kubectl OOM remediations on a latency alert).
\texttt{1a95f6a6} has placeholder evidence.  \texttt{79119df0} has
zero evidence and zero actions.  \texttt{d763a9c9} cites a
truncation artefact as cause.  \texttt{0c76258b} has placeholder
actions.  \texttt{be7019bf} has the explicit
\texttt{Shelved --- awaiting recurrence} non-action.
\texttt{8891233d} names no cause for a 99.6\%-memory pod.  Sprint~4
P11 (\texttt{S4-PR-04}) tightens the classifier and P8
(\texttt{S4-PR-05}) reclassifies the corpus.

\paragraph{F4.~Diagnostic-step boilerplate (P3).}
Eight of twelve sampled RCAs that have non-empty
\texttt{diagnostic\_steps} draw from the same four-string
clamp-fallback template (\texttt{Open Grafana Explore\ldots},
\texttt{Check kube events\ldots}, \texttt{Open Jaeger and inspect
the slowest trace\ldots}, \texttt{Do NOT run kubectl
rollout/scale/set\ldots}).  Two distinct \texttt{TargetDown} fires
(\texttt{64e46e6c}, \texttt{a9393ef6}) and two distinct
\texttt{PodHighMemoryUsage} fires (\texttt{0de9801c},
\texttt{8891233d}) carry near-identical step sets.  Sprint~4 P3
(\texttt{S4-PR-02}) moves generation into a deterministic
templater fed from concrete evidence.

\paragraph{F5.~Wrong-archetype suggested actions (P8).}
\texttt{0b215ef3} suggests \texttt{kubectl set resources
deploy/spring-boot --limits=memory=2Gi} for a Kong latency alert;
\texttt{d763a9c9} suggests memory-limit + JVM-flag changes for a
truncated-metric-name hallucination.  Both rows sit in the
exemplar corpus as \texttt{actionable} and thus pollute future
RCAs that retrieve them via \texttt{get\_similar\_decisions}.
Sprint~4 P8 (\texttt{S4-PR-05}) manually demotes
\texttt{0b215ef3} and adds a nightly post-hoc revalidator.

\paragraph{F6.~No-action escalations (P5).}
\texttt{79119df0}, \texttt{be7019bf}, \texttt{8891233d} all carry
\texttt{verdict=escalate} (or \texttt{verdict=dismiss} with a
\texttt{shelved} non-action) with no concrete suggested action.
The 2026-05-20
\texttt{PodHighMemoryUsage} case (\texttt{8891233d}) is the
specific incident that produced the Sprint-4 carry story
\texttt{S4-HF-01} (widen bounded-agency trigger to low-confidence
+ clamp-fired verdicts): the system saw memory at 99.6\% and
returned five debug-log lines, no named cause, no actions, but
verdict \texttt{escalate}.

\subsection{What is working}

The audit also confirms several mechanisms behaving exactly as
designed:

\begin{itemize}
  \item \textbf{Adaptive thresholds with $\sigma$-claims} (US-5.4)
        in case \texttt{98bef564}: ``current CPU usage of 3.4\% is
        +0.4$\sigma$ above the same-hour baseline from one week
        ago (\texttt{stddev\_over\_time} + offset 7d), well within
        the natural seasonality of the load curve.''  The rule-design
        framing is correct and actionable.
  \item \textbf{Synthetic-alert detection} in case
        \texttt{c29e557b}: the system identifies an
        \texttt{audit-live.sh}-fired
        \texttt{HighP95Latency} probe and correctly labels it
        \texttt{data\_starved} with \texttt{verdict=dismiss}.
  \item \textbf{Deep-trace integration} (S3-HF-07) in case
        \texttt{0de9801c}: the RCA cites
        \texttt{Jaeger trace 7f3a2c1d9b4e: GET /api/employee took
        2347 ms, span waits on MySQL connection pool} ---
        a concrete trace ID with span-level evidence.  This is the
        bounded-agency MCP path working as intended.
  \item \textbf{Transient-fault discrimination} in case
        \texttt{64e46e6c}: a two-scrape-interval \texttt{up=0}
        flap is correctly identified as ``transient scrape
        collision or network jitter, not an actual outage''.  The
        absence of corroborating Drain3 or host-metric signal is
        used as the discriminating evidence, not just dismissed
        as silence.
\end{itemize}

\subsection{Overall verdict and Sprint-4 alignment}

Three of the twelve sampled RCAs would survive a strict operator
review unmodified (\texttt{64e46e6c}, \texttt{98bef564},
\texttt{c29e557b}); two more are borderline-acceptable
(\texttt{a9393ef6}, \texttt{0de9801c}); seven exhibit at least one
structural defect that the Sprint-4 P0--P11 plan was specifically
scoped to fix.  The distribution matches expectations: the
structural foundation (MCP-only invariant, exemplar retrieval,
confidence clamp, recurrence gate, deep-trace path) is sound, and
the binding constraint on output quality lives one layer up --- in
the prompt-assembly thinness (P1, F2), the exemplar-substitution
gap (P6, F1), the over-permissive quality classifier (P11, F3),
and the diagnostic-step layer placement (P3, F4).  The same
binding constraint was independently identified by the
2026-05-19 heavy-load investigation that produced the P0--P11
plan; the critical audit confirms the priorities without
revising them.

The next critical re-audit is scheduled for the close of Sprint~4,
against the same 8-criterion rubric.  Sprint-4 acceptance: the
pass rate for C4 (no-noise) is expected to move from 33\% to
$\geq$80\%, C7 (\texttt{rca\_quality} accuracy) from 33\% to
$\geq$70\%, C6 (non-templated diagnostic steps) from 17\% to
$\geq$60\%.  If those three move, the structural-foundation /
output-quality framing of this audit is empirically validated.

\section{Chaos-run outcomes}

The chaos-test harness in
\texttt{monitoring-project/scripts/chaos/} runs five scenarios
(\texttt{target\_down}, \texttt{high\_memory\_usage},
\texttt{high\_cpu\_usage}, \texttt{drain3\_anomaly}, and the new
\texttt{0b215ef3}-style corpus seed).  Each run induces a controlled
failure, polls \texttt{/decisions} for a matching alert, scores the
RCA, and regenerates \texttt{chaos-report-YYYY-MM-DD.html}.

Run 3 (2026-04-28) captured 0/4 because pod-scoped chaos did not
move host-level metrics.  This finding produced the pod-level alert
rules (\texttt{PodHighMemoryUsage}, \texttt{PodHighCpuUsage}) added
the same evening.  Run 4 (2026-04-28 PM) captured 0/2, this time
for known reasons: container restart churn split timeseries across
\texttt{cgroup} identifiers; cumulative CPU expressions did not match
their denominator labels.  Both are tractable thirty-minute fixes
filed as Sprint~4 candidate~\#16.

\section{Audit-corpus consistency}

The daily static audit ran 17 times during the investigation window
(2026-04-30 to 2026-05-19), all green at 226/226 with the same
seven-item doc-drift list carried forward day-on-day --- no new drift
accumulated.  At sprint close-out (2026-05-19) the paper-cut sweep
closed all seven items and the suite stepped to \texttt{252/252}
following the Epic-4 close-out commits.  The live audit (cron on
\texttt{claude-controller}) silently stopped firing after 2026-04-29
due to a controller-uptime issue; the cron was reinstalled on
2026-05-19 and is now firing daily.  The 20-day gap is filed as
resilience item R-1 (chapter~\ref{ch:future}).

\section{Latency budget}

The Sprint~2 verdict-story field note recorded an end-to-end pipeline
latency of \texttt{67\,995\,ms} ($\approx$\,68\,s) on the laptop
runner with \texttt{qwen2.5:7b-instruct}, against a historical
k3s-on-CPU baseline of 20--40 \emph{minutes} per alert.  Sustained
steady-state latencies on the laptop are 25--35\,s warm.  Five
recent fires were anomalously slow (323, 399, 533, 755, 1081\,s)
and the per-stage instrumentation drafted during the investigation
window will surface whether the cause is bounded-agency retries,
LLM prefill saturation, or MCP tail latency; the instrumentation is
held until S3-HF-05 ships so the retry path is properly attributable.

\section{Operational steady state}

The qualitative observation that the platform held a steady state
across the twenty-day investigation window --- plus a subsequent
sprint close-out wave on 2026-05-19 that landed ten SP of refactors
across four repositories without a regression --- is itself an
evaluation result.  No incidents were observed across the
twenty-seven-day run from the Sprint-2 close-out on 2026-04-23 to the
verification on 2026-05-20; AWS instances stayed up across the full
window; the \texttt{ai-triage-service} container on the laptop
required exactly one manual restart over the period (caused by
Docker restart-backoff after a WSL sleep, not by a platform fault).
This validates the architectural choice of running the AI layer
separately from the production infrastructure: a transient laptop
fault does not affect the observability stack, only the AI-assisted
RCAs.

% =====================================================================
\chapter{Design Decisions and Rationale}
\label{ch:decisions}

The project maintains a durable decisions log
(\texttt{monitoring-docs/decisions-log.html}) recording the
twenty-three substantive design or rule-of-thumb decisions taken
during construction.  Each entry follows a stable structure: the
problem observed, the constraints, the decision, and the verification
that closed it.  Table~\ref{tab:decisions} summarises the catalogue.

\begin{longtable}{p{0.04\textwidth} p{0.34\textwidth} p{0.56\textwidth}}
\caption{Catalogue of durable design decisions D1--D23.}
\label{tab:decisions} \\
\toprule
\textbf{ID} & \textbf{Title} & \textbf{One-line decision} \\
\midrule
\endfirsthead
\toprule
\textbf{ID} & \textbf{Title} & \textbf{One-line decision} \\
\midrule
\endhead
\midrule
\multicolumn{3}{r}{\itshape (continued)} \\
\endfoot
\bottomrule
\endlastfoot
D1  & LLM hedging with ``insufficient data''
    & Treat ``insufficient data'' as a verdict, not a substitute for
      one; persist it but never escalate on it. \\
D2  & Empty \texttt{suggested\_actions} on most RCAs
    & Add the action-template fallback layer; never emit an RCA with
      both empty actions and \texttt{rca\_quality=actionable}. \\
D3  & Observed value picked from the wrong \texttt{refId}
    & Default to refId=A; require an explicit override in the alert
      rule to use a different reduced expression. \\
D4  & \texttt{llm\_confidence} never persisted
    & Persist confidence; do not surface a row without it; use the
      value to gate the action-template fallback. \\
D5  & \texttt{rca\_quality} lied when actions+evidence empty
    & Compute quality from evidence and action shape, not from the
      verdict string. \\
D6  & Overnight email incident --- rogue k3s AI stack
    & Decommission k3s-hosted AI; move AI onto laptop;
      remove Alertmanager; unify alerting in Grafana. \\
D7  & Deployment-type mismatch
    & Carry an explicit \texttt{service\_deployment\_type} map and
      template actions by deployment type. \\
D8  & Drain3 anomalies never fired into triage
    & Wire a synthetic webhook from the Drain3 poller into
      \texttt{/webhook/drain3}; relax thresholds (rate 0.10 $\to$
      0.02, lines 100 $\to$ 20). \\
D9  & Flapping alerts consumed $N\times$ LLM calls
    & Two-tier dedup; pre-LLM fingerprint suppression; recurrence
      gates (later, D21). \\
D10 & Correlated alerts rendered but not prompted
    & Include the correlated-alerts list in the prompt's
      pre-gathered context. \\
D11 & Monitoring stack stays on Docker Compose
    & Migrate the application workload to k3s; keep observability on
      Docker Compose. \\
D12 & Structured outputs vs.\ JSON mode
    & Use structured-output schema enforcement at the LLM layer;
      retry on schema parse failure. \\
D13 & Pre-LLM metric interpreter
    & Move numeric reasoning out of the LLM into a deterministic
      Python module; pass interpreted values into the prompt. \\
D14 & Downtime backfill from Grafana annotations
    & At startup, query Grafana annotations for the missing window,
      replay them through the pipeline (bounded by a max gap). \\
D15 & Bounded-agency retry
    & Exactly one MCP call from a six-tool whitelist; revalidate;
      keep first-pass if retry still bad. \\
D16 & UEBA framing for Sprint 2 Epic 5
    & Add the entity-behavior layer (per-service Drain3, baselines,
      feedback) as an extension to the event-centric core. \\
D17 & Exemplar library
    & Inject one of fourteen canonical RCA archetypes as structural
      scaffolding before the evidence pack. \\
D18 & Adaptive thresholds
    & Replace static numeric thresholds on the three flappiest rules
      with same-hour-7-days-ago + $3\sigma$ baselines. \\
D19 & RCA prose quality
    & Telemetry is the means, not the end --- RCAs must name a
      cause, not restate the alert. \\
D20 & Closed-loop feedback
    & Capture operator overrides; compute precision/recall;
      consume overrides in the suppression layer. \\
D21 & Recurrence gates
    & Pre-LLM + post-LLM hysteresis for opted-in alerts; critical
      severity bypasses. \\
D22 & SESSION\_HANDOFF kept controller-local
    & Untrack from public repo; scrub history;
      keep as a local scratchpad. \\
D23 & Per-service Drain3 template trees
    & Dictionary of \texttt{service $\to$ TemplateMiner};
      file persistence per service; migration of pre-split state. \\
\end{longtable}

\section{Three decisions in depth}

\subsection{D17 --- Why exemplars work where retrieval did not}

The naive approach to ``teach the LLM what good looks like'' would
have been retrieval-augmented generation over the
\texttt{rca\_history.db} table.  D17 rejects this for two reasons.
First, the historical record \emph{contains the failures we are
trying to fix} --- training the model on the empirical RCA history
would reinforce the surface-only and hypothesis-menu patterns the
project is trying to eliminate.  Second, the historical record is
unlabelled at the time of the decision; without operator-feedback
quality ratings, retrieval would pick by surface similarity, not by
quality.

The exemplar library is a curated corpus, not an empirical one.  Each
of the fourteen archetypes is hand-authored by the engineering team
to model the prose shape and reasoning structure the platform wants
to reproduce.  This is a deliberate inversion of the usual
``train on what we have'' approach in favour of ``train on what we
want''.  The BM25 retrieval dry-run during the investigation window
(\S\ref{sec:investigation}) confirmed empirically that retrieval over
fourteen archetypes is degenerate: the top-1 hit always matches the
existing \texttt{find\_for\_alert} selector, so adding retrieval on
top contributes nothing.  Once the operator-feedback table accumulates
quality-rated entries (Epic 2 work), curated retrieval over a
quality-filtered subset becomes worthwhile --- this is the Sprint~4
US-3.3 / US-3.4 scope.

\subsection{The MCP-only invariant as a project-wide rule}

The MCP-only invariant is not a technical mechanism but a
project-wide rule.  Its enforcement is partly social (code review,
the audit reports), partly architectural (the MCP servers exist;
they have error envelopes and metrics), and ultimately mechanical
(the S3-HF-08 continuous-integration lint forbids the bypass
patterns).  The \texttt{0b215ef3} incident demonstrated the cost of
relying on the social layer alone: three of the project's documents
(\texttt{prepare.html}, \texttt{system-explained.html},
\texttt{triage-service.html}) claimed an invariant the code did not
enforce, and the resulting RCA carried hallucinated content because
of it.

The mechanism that makes the rule durable is the combination of
(a) error-envelope-shaped MCP responses, so a failed query is
observable rather than silent; (b) MCP-bridge metrics, so call rates
and latencies are visible in Grafana; and (c) the CI lint, so the
rule has teeth on new code as well as existing.  The order of
operations matters: the lint cannot fire until S3-HF-04 and S3-HF-05
have closed the existing offenders, because a lint that fails on
existing code is one the team will work around rather than fix.

\subsection{Bounded agency over full agentic gather}

Decision~D15 introduces the bounded-agency retry: the LLM is given
exactly one extra MCP call after a validator violation, from a
six-tool whitelist, with a retry-feedback string and the same
Pydantic schema as the first pass.  This is a deliberately constrained
form of agentic gather: an unbounded loop (Sprint-4 candidate \#12,
Tier~3) gives better RCA quality on data-starved cases but introduces
unbounded latency and unbounded tool-call cost.  Bounded agency is
the project's currently-chosen trade-off: predictable latency, bounded
LLM cost per alert, and a measurable quality improvement on
validator-flagged cases.

The natural extension is Tier~3 (iterative agentic gather with
hypothesis-tree state); the project scopes it to Sprint~4 explicitly
gated on Tier~4 (the chaos-scorer regression gate from S3-HF-09).
Shipping Tier~3 without Tier~4 in place would be a high-risk change
without a measurement scaffold; this gating is a deliberate
methodological choice.

% =====================================================================
\chapter{Future Work and Sprint 4 Roadmap}
\label{ch:future}

\section{Sprint 4 candidate ranking}

The ranked list of Sprint~4 candidates as of 2026-05-19 is reproduced
in table~\ref{tab:sprint4}.  The list integrates the twelve findings
(P0--P11) from the 2026-05-19 real-load investigation with the
existing carry-forward backlog items, ranked by
$(\text{impact}\times\text{frequency})\div\text{effort}$.  The two
anchor items are \textbf{(P0+P9)} \emph{Latency bundle} --- the
latency unblocker, attacked as a bundle of small fixes after the
GPU-migration option was dropped on cost grounds (\S\ref{sec:gpu-dropped}) ---
and \textbf{P2} \emph{Suppression-cascade fixes} --- the
cheapest-impact win, $\sim$1.5\,h for a class of false-negative
caught twice in a single day's evidence.

\begin{longtable}{rp{0.40\textwidth} p{0.09\textwidth} p{0.07\textwidth} p{0.22\textwidth}}
\caption{Sprint 4 candidates as of 2026-05-19, ranked.  ``P-tag''
maps each row to the 2026-05-19 investigation findings P0--P11
where applicable; \texttt{---} indicates a carry-forward item with
no investigation finding attached.}
\label{tab:sprint4} \\
\toprule
\textbf{\#} & \textbf{Item} & \textbf{Effort} & \textbf{P-tag} & \textbf{Gates} \\
\midrule
\endfirsthead
\toprule
\textbf{\#} & \textbf{Item} & \textbf{Effort} & \textbf{P-tag} & \textbf{Gates} \\
\midrule
\endhead
\bottomrule
\endfoot
1 & \textbf{Latency bundle} --- hard 180\,s pipeline timeout +
    per-stage instrumentation + Ollama concurrency cap (serial).
    Combined target: median 6\,min $\to$ $\sim$90\,s on laptop.
                                              & $\sim$5\,h total
                                              & \textbf{P0+P9}
                                              & None \\
2 & \textbf{Suppression-cascade fixes} --- exclude
    \texttt{audit-live-*}/\texttt{chaos-*} fingerprints from
    suppression lookup; tighten
    \texttt{TRIAGE\_HISTORY\_LOOKBACK\_MINUTES} 30\,$\to$\,10;
    widen recurrence-gate annotation to also cover
    \texttt{TargetDown}, \texttt{HighP95Latency},
    \texttt{HighKongP95Latency}.
                                              & $\sim$1.5\,h
                                              & \textbf{P2}
                                              & None --- cheapest impact \\
3 & \textbf{PromQL / numerical-noise scrubber} --- validator rule
    rejecting PromQL function calls, raw floats $>$\,4 sig figs,
    Z-score equations, and \texttt{...}-truncated metric names;
    \texttt{translate\_value(value,unit)} helper renders values
    before prompt assembly.
                                              & $\sim$4\,h
                                              & \textbf{P1}
                                              & None \\
4 & \textbf{Deterministic diagnostic-steps generator} --- move
    \texttt{diagnostic\_steps} out of the LLM into
    \texttt{app/diagnostic\_steps\_for\_alert(alert,ctx)}, templated
    from concrete evidence (trace\_id, service, window, observed
    value).
                                              & $\sim$6\,h
                                              & \textbf{P3}
                                              & None \\
5 & \textbf{US-5.2 Incident correlator} --- collapse a 4-alert
    kill-chain RCA into one coherent email.
                                              & 1--2\,d
                                              & \textbf{P10}
                                              & None \\
6 & \textbf{Drain3 + empty-actions auto-downgrade} --- post-LLM
    gate: if \texttt{verdict=escalate} \textsc{and}
    \texttt{suggested\_actions=[]} \textsc{and} no service named,
    auto-downgrade to dismiss.  Drain3 escalations require
    $\ge$\,2 novel templates in a 5-min window.
                                              & 2--3\,h
                                              & \textbf{P5}
                                              & None \\
7 & \textbf{Hallucination on truncated metric names} --- context
    sanitiser drops lines ending in \texttt{(truncated)} or
    \texttt{...]}; cross-check cited metric names against
    \texttt{app/metrics.py} and the Prometheus registry.
                                              & $\sim$4\,h
                                              & \textbf{P4}
                                              & None \\
8 & \textbf{Tighter \texttt{\_classify\_rca\_quality}} ---
    actionable requires ($\ge$\,1 specific evidence item)
    \textsc{and} (named cause in prose) \textsc{and} ($\ge$\,1
    action \textsc{or} explicit ``dismiss with reasoning'').
                                              & 1\,h
                                              & \textbf{P11}
                                              & None \\
9 & \textbf{Corpus reclassification + nightly post-hoc} ---
    manually demote \texttt{0b215ef3} and clear its actions; a
    nightly job re-runs the current validator over the last 7\,d of
    stored RCAs, demoting any that would now fail.
                                              & 4\,h
                                              & \textbf{P8}
                                              & P11 first \\
10 & \textbf{MCP native tool-calling rewrite} --- replaces
     prompt-engineered JSON-schema output with the native tool-call
     API; smaller prompts $\to$ faster inference, compounds with
     item \#1.
                                              & 5\,SP
                                              & ---
                                              & GPU gate removed \\
11 & \textbf{Exemplar prose rewrites + overlap-detection
     validator} --- exemplars use \texttt{\{INSTANCE\}},
     \texttt{\{N\}}, \texttt{\{TIMESTAMP\}} placeholders the LLM
     must substitute; validator rejects RCAs with $>$\,80\,\%
     char-overlap to their matched exemplar.
                                              & $\sim$1\,d
                                              & \textbf{P6}
                                              & None \\
12 & \textbf{Tier 3 --- Iterative agentic gather} (3-iteration
     tool-call loop with hypothesis-tree state)
                                              & 3--5\,d
                                              & ---
                                              & S3-HF-09 $\checkmark$ \\
13 & \textbf{Curated RAG retrieval} (BM25 over feedback-curated
     YAML; top-$k$\,=\,3 positives + 2 anti-examples)
                                              & was US-3.4
                                              & ---
                                              & Epic 2 feedback table populated \\
14 & \textbf{Weekly curation job} --- APScheduler, Sundays 02:00
     UTC, reads SQLite feedback into \texttt{library\_curated.yaml}
                                              & was US-3.3
                                              & ---
                                              & Feedback volume \\
15 & \textbf{Dashboard \texttt{dismiss\_with\_recommendation}
     colour} --- amber for ``dismiss but rule needs fixing'' vs.\
     grey for ``dismiss, ignore.''  Schema migration on
     \texttt{rca\_history}.
                                              & 2\,h
                                              & \textbf{P7}
                                              & P11 first \\
16 & \textbf{Pod-level alert rule fixes} --- memory cgroup-churn
     aggregation + CPU label mismatch.
                                              & $\sim$30\,min each
                                              & ---
                                              & None \\
17 & \textbf{O-9 webhook auth} + \textbf{O-10 nightly RCA-history
     S3 backup} --- Friday-slot hardening.
                                              & $\sim$30\,min each
                                              & ---
                                              & None \\
18 & Sentence-transformers retriever (RAG fallback) & ---
                                              & ---
                                              & Only if BM25 underperforms \\
19 & Trace-ID linkage from log anomalies     & 2\,SP
                                              & ---
                                              & None \\
20 & RCA playbook templates                   & ---
                                              & ---
                                              & Complements Epic 2 \\
21 & US-5.7 labelled corpus expansion + F1 tracking
                                              & ---
                                              & ---
                                              & S3-HF-09 $\checkmark$ \\
22 & L2/L3 chaos scenarios (O-12)             & ---
                                              & ---
                                              & None \\
23 & Drain3 self-monitoring loop allowlist
     (\texttt{service\_name=drain3} exclusion)
                                              & $\sim$30\,min
                                              & ---
                                              & None \\
24 & Doc debt --- \texttt{monitoring-project/CLAUDE.md} refresh
     + stale architecture PNGs replaced.
                                              & 1--2\,h
                                              & ---
                                              & None \\
\end{longtable}

\noindent
\textbf{Suggested week-1 execution order:}
\#1 $\to$ \#2 $\to$ \#3 $\to$ \#4\,+\,\#8 (paired).  Items
\#1\,+\,\#2\,+\,\#8 together are roughly eight hours of effort and
deliver the largest perceived-quality improvement.  Items \#3\,+\,\#4
are the ``the prose is now honest'' work.  Items \#5--\#11 close the
false-positive escalation paths.  Items \#12--\#24 are wave 2.

\section{Tier 3 --- iterative agentic gather}

The Tier-3 design replaces the single-shot bounded-agency retry with
a three-iteration tool-call loop, each iteration producing a partial
hypothesis tree that the model expands on the next iteration.  The
hypothesis tree is stored in the prompt as a structured section; at
each iteration the model is asked to expand the most promising leaf,
gather supporting evidence via a single MCP call, and update the
tree.  Termination conditions: a confidence threshold reached, a
maximum-depth limit hit, or a total latency budget exhausted.

This is materially more permissive than the bounded-agency retry;
the gating on S3-HF-09 reflects the project's view that shipping
Tier~3 without an automated quality-regression gate would be
imprudent.

\section{Curated RAG retrieval (US-3.3 / US-3.4)}

Once the Epic-2 operator feedback table has accumulated rated entries,
curated retrieval becomes worth the engineering effort.  The design:
a weekly job reads the feedback table for the last seven days; entries
above a quality threshold are written as YAML to a
\texttt{library\_curated.yaml} file; the triage prompt's exemplar
section is supplemented by a small number ($k=3$ positives, $k=2$
anti-examples) retrieved by BM25 over that curated set.  The
investigation-window dry-run established that BM25 over the
fourteen-archetype \emph{base} library is degenerate but did not
evaluate BM25 over an enriched library; that is the actual
production scenario and will be measured once Epic~2 ships.

\section{Resilience tier (R-1 to R-5)}

The 20-day investigation window surfaced five operational-resilience
items that were not blocking but are worth Sprint~4 inclusion:

\begin{itemize}
  \item \textbf{R-1} Live-cron heartbeat metric on a Prometheus
        push-gateway with a \texttt{LiveAuditMissing} Grafana rule.
  \item \textbf{R-2} Move the live-audit cron off the ephemeral
        \texttt{claude-controller} onto the durable
        \texttt{observability-rca-monitoring} EC2 via a
        \texttt{systemd} timer.
  \item \textbf{R-3} CI lint that fails if any
        \texttt{sprint*-plan.md} window end-date is in the past.
  \item \textbf{R-4} \texttt{the-bigger-picture.md} staleness badge in
        the repository README, red at fourteen days.
  \item \textbf{R-5} \texttt{SESSION\_HANDOFF.md} S3 backup, folded
        into the O-10 nightly-backup work.
\end{itemize}

\section{Dropped scope: cloud GPU migration}
\label{sec:gpu-dropped}

The cloud-GPU migration to \texttt{us-west-2} (\texttt{g5.xlarge}
A10G 24\,GB, host for \texttt{qwen2.5:14b} or \texttt{32b-q4}) sat as
a Sprint-4 candidate from 2026-04-27, when the
\texttt{G+VT} on-demand vCPU quota was approved.  It is
\textbf{dropped on cost grounds on 2026-05-19}.

A \texttt{g5.xlarge} runs approximately one US dollar per hour
on-demand, or roughly \$720 per month if left running, with
aggressive auto-stop bringing the realistic monthly cost to
\$200--300.  This is an order of magnitude more than the laptop
GTX~1060 + \texttt{qwen2.5:7b} setup that has been the production
runner for twenty-seven days with no quality regression surfaced in
the daily audits or the investigation-window RCA scoring.  The
\texttt{qwen2.5:14b} / \texttt{32b-q4} upgrade path is parked
indefinitely; the gating condition for revisiting the decision is
either a material degradation of 7b output quality or the
availability of a free or heavily subsidised GPU.

The decision has two knock-on effects.  First, the MCP native
tool-calling rewrite (Sprint~4 candidate~\#10) was previously gated
on the GPU migration; the larger model is more reliable on
function-calling benchmarks than the 7B and the rewrite was held
until the migration shipped.  With the migration dropped, the
rewrite proceeds on the 7B model directly; the structured-output
schema enforcement and the bounded-agency retry combine to keep
tool-calling reliability operationally adequate at this scale.
Second, the \texttt{us-west-2} G+VT quota approval (4 vCPU, granted
2026-04-27) remains on the account but stays unused.

\section{Beyond Sprint 4}

Beyond the ranked candidates, three longer-arc directions are worth
recording for the supervisor's consideration:

\begin{itemize}
  \item A \textbf{multi-tenant deployment model}.  The current
        platform is single-tenant by construction.  A multi-tenant
        deployment (e.g.\ one triage service per business unit at
        Tanger Med, with shared MCP bridges and a shared
        observability stack) would require namespace isolation in the
        RCA history, per-tenant exemplar libraries, and per-tenant
        confidence calibration.
  \item \textbf{Direct integration with Tanger Med production
        systems}.  The platform is currently exercised against a
        representative web stack.  Onboarding actual production
        services would require careful PII handling (operator
        feedback content, log lines), an SLO programme to set
        meaningful confidence thresholds, and a formal change-control
        process for prompt and exemplar updates.
  \item \textbf{Federated model evaluation}.  As the project
        accumulates labelled RCA outcomes via the feedback layer,
        comparing alternative LLM choices on a held-out set becomes
        feasible.  Models worth evaluating: \texttt{qwen2.5:14b-instruct},
        \texttt{deepseek-r1:14b-distill}, \texttt{mistral-nemo:12b}.
        The cost-vs-quality curve will be the deciding factor; this
        is the natural successor to the dropped GPU migration.
\end{itemize}

% =====================================================================
\chapter{Conclusion}
\label{ch:conclusion}

This dissertation has reported on the design and implementation of
the CIRES Observability Platform, an AI-augmented observability
solution built over three completed sprints at CIRES Technologies.
The platform addresses the operational gap identified during Sprint~0
field research --- that the team was finding out about production
problems after the fact rather than proactively --- by combining a
conventional open-source observability stack with an AI triage layer
that turns each alert into a structured, pre-investigated RCA.

The principal contributions of the work are:

\begin{enumerate}
  \item \textbf{An end-to-end reference implementation.}  A
        Terraform-provisioned, Ansible-configured, k3s-and-Docker
        observability stack on AWS, with a FastAPI-based AI triage
        service producing structured RCAs from Grafana webhooks via
        a locally-hosted LLM.  Reproducible from clean infrastructure
        and operable from a single playbook command.
  \item \textbf{The MCP-only data-access invariant.}  A project-wide
        rule, partly architectural and partly mechanical (via the
        Sprint-4 lint), that every datum the LLM sees must arrive
        through an MCP bridge.  The rule produces a provenance
        regime over the LLM's input and is enforced via the bridge
        layer's metrics, the audit reports, and a continuous-integration
        lint.
  \item \textbf{The hallucination firewall.}  A layered combination
        of mechanisms (exemplar injection, alert-keyed blocklist,
        surface-only validator, hypothesis-menu validator,
        cause-evidence rule, confidence clamp with action stripping,
        bounded-agency retry) designed to keep LLM output
        operationally trustworthy.  Each mechanism addresses a
        specific failure mode that was observed empirically during
        construction and is recorded in the decisions log with the
        evidence that drove it.
  \item \textbf{The closed-loop operator-feedback layer.}  A schema
        and surface that captures operator overrides on RCAs and
        propagates them into the suppression layer, the precision/recall
        metrics, and (in Sprint~4) the curated retrieval corpus.
  \item \textbf{A documented decisions log.}  Twenty-three durable
        design decisions (D1--D23) recorded contemporaneously with
        the engineering work, including the problem observed, the
        constraints, the decision, and the verification.  The log is
        a primary deliverable of the project and is exercisable in
        future audits.
\end{enumerate}

The platform is in active service against a representative web stack
on AWS, with daily static and live audits producing a documented
trail of operational health.  The remaining Sprint~3 work and the
ranked Sprint~4 candidates carry the platform forward; the
methodological choice to gate Sprint~4's iterative agentic gather
(Tier 3) on a measurement scaffold (Tier 4) is, in the author's view,
the most important durable lesson from the project --- the cost of
shipping an unmeasurable AI mechanism into production materially
exceeds the cost of building the measurement first.

% =====================================================================
\appendix
\chapter{Repository inventory}
\label{app:repos}

\begin{table}[h]
\centering
\small
\begin{tabular}{ll}
\toprule
\textbf{Repository} & \textbf{Purpose} \\
\midrule
\texttt{monitoring-project}            & Ansible roles, playbooks,
                                         chaos harness \\
\texttt{provisioning-monitoring-infra} & Terraform AWS state \\
\texttt{monitoring-docs}               & This document, decisions
                                         log, architecture page,
                                         sprint history \\
\texttt{monitoring-triage-service}     & FastAPI triage service,
                                         pipeline, validators \\
\texttt{monitoring-mcp-servers}        & Five MCP bridges \\
\texttt{monitoring-audit-results}      & Private --- daily static +
                                         live audits \\
\texttt{jira}                          & Jira CSV exports and import
                                         scripts \\
\bottomrule
\end{tabular}
\caption{Project repository layout (seven repositories).}
\label{tab:repos}
\end{table}

\chapter{Glossary}
\label{app:glossary}

\begin{description}
\item[AIOps] An informal label for the application of artificial-
    intelligence techniques to IT operations problems --- anomaly
    detection, root-cause analysis, alert summarisation.
\item[Bounded-agency retry] The platform's one-MCP-call retry layer.
    The LLM is given exactly one additional tool call after a
    validator violation; the schema and the retry-feedback string
    constrain the call.
\item[Cause-evidence rule] A validator rule that tokenises the RCA
    prose and the cited evidence with \texttt{[a-z]\{4,\}} and
    requires a non-stop-word overlap.  Rejects RCAs whose cause is
    not supported by the cited evidence.
\item[Confidence clamp] The pipeline step that forces
    \texttt{confidence} $\leq 0.4$ if the surface-only or
    hypothesis-menu validators fired, if the RCA quality is
    \texttt{data\_starved}, or if the actions look templated.  Strips
    \texttt{suggested\_actions} when fired and emits read-only
    \texttt{diagnostic\_steps} instead.
\item[Drain3] A log-template clustering algorithm (a tunable variant
    of Drain).  Used in the platform as a fourth telemetry pillar
    surfacing novel templates.
\item[Exemplar archetype] One of fourteen hand-authored RCA examples
    injected into the prompt as structural scaffolding before the
    evidence pack.
\item[Grafana unified alerting] The successor to Alertmanager in the
    Grafana stack.  All alert rules and contact points live inside
    Grafana itself.
\item[Hallucination] An LLM output statement that is not entailed by
    the input context (e.g.\ a service that does not exist; a
    metric value that was not in the evidence pack).
\item[Hypothesis menu] An RCA that hedges with multiple possible
    causes.  Detected by the S3-HF-03 validator.
\item[MCP] Model-Context-Protocol --- an open specification for
    connecting LLM applications to external data sources and tools.
\item[MCP-only invariant] The project-wide rule that every datum the
    LLM sees must arrive through an MCP server.
\item[node-exporter] The standard Prometheus exporter for Unix host
    metrics (CPU, memory, disk, filesystem, network).
\item[OTel] OpenTelemetry, the de-facto open standard for telemetry
    instrumentation.
\item[RCA] Root-Cause Analysis --- the structured explanation of why
    an incident occurred.
\item[SRE] Site Reliability Engineering --- the operational discipline
    formalised by Google's homonymous book and workbook.
\item[Surface-only LEDE] An RCA opening that restates the alert (e.g.\
    leading with the PromQL expression or the observed value)
    instead of naming a cause.  Detected by the
    \texttt{\_SURFACE\_ONLY\_LEDE\_PATTERNS} regex set.
\item[Tailscale tailnet] The WireGuard-based mesh that connects all
    hosts in the platform.  MagicDNS provides stable host names.
\item[Triage] In the platform's vocabulary, the AI layer that consumes
    alerts and produces RCAs (as opposed to the human-facing
    triage practice in incident response).
\end{description}

\chapter{Sprint chronology summary}
\label{app:sprint-summary}

\begin{table}[h]
\centering
\small
\begin{tabular}{p{0.07\textwidth} p{0.30\textwidth} p{0.55\textwidth}}
\toprule
\textbf{Sprint} & \textbf{Window} & \textbf{Headline outcome} \\
\midrule
0 & Feb 2026 -- early Mar 2026 & Problem statement crystallised; solution
hypothesis (observability + AI triage) chosen. \\
1 & 2026-03-10 -- 2026-03-30 & On-prem three-VM observability stack
deployable in 15\,min via Ansible. \\
2 & 2026-03-31 -- 2026-04-23 & AWS migration, k3s for the application
workload, AI RCA MVP (5 MCP servers + FastAPI triage + Ollama). \\
2-ext & 2026-04-23 -- 2026-04-29 & Epic 5: per-service Drain3,
entity baselines, closed-loop feedback, recurrence gates, adaptive
thresholds. \\
3 & 2026-04-23 -- 2026-05-19 & RCA quality close-out, hallucination
firewall epic added on 2026-04-29 after 0b215ef3, 20-day investigation
window 04-30 -- 05-19. \\
4 & TBD & Ranked candidate list: incident correlator, Tier 3 agentic
gather, curated RAG, pod-rule fixes, R-1..R-5 resilience tier. GPU
migration dropped on cost grounds. \\
\bottomrule
\end{tabular}
\caption{Sprint chronology, condensed.}
\label{tab:sprint-summary}
\end{table}

% ---------- bibliography ----------
\chapter*{References}
\addcontentsline{toc}{chapter}{References}
\markboth{References}{References}

\begin{thebibliography}{99}

\bibitem{sre-book}
B.~Beyer, C.~Jones, J.~Petoff, N.~R.~Murphy (eds.),
\emph{Site Reliability Engineering: How Google Runs Production Systems},
O'Reilly, 2016.

\bibitem{sre-workbook}
B.~Beyer, N.~R.~Murphy, D.~K.~Rensin, K.~Kawahara, S.~Thorne (eds.),
\emph{The Site Reliability Workbook: Practical Ways to Implement SRE},
O'Reilly, 2018.

\bibitem{drain}
P.~He, J.~Zhu, Z.~Zheng, M.~R.~Lyu, \emph{Drain: An Online Log Parsing
Approach with Fixed Depth Tree}, IEEE International Conference on Web
Services (ICWS), 2017.

\bibitem{otel}
OpenTelemetry Project,
\emph{OpenTelemetry Specification}, version 1.x,
\url{https://opentelemetry.io/docs/specs/otel/}.

\bibitem{prometheus}
J.~Volz et al.,
\emph{Prometheus: Up \& Running}, O'Reilly, 2018, and
\url{https://prometheus.io/docs/}.

\bibitem{loki}
Grafana Labs,
\emph{Grafana Loki documentation}, \url{https://grafana.com/docs/loki/latest/}.

\bibitem{jaeger}
Y.~Shkuro,
\emph{Mastering Distributed Tracing}, Packt, 2019,
and \url{https://www.jaegertracing.io/docs/}.

\bibitem{mcp}
Anthropic,
\emph{Model Context Protocol --- specification},
\url{https://modelcontextprotocol.io/}.

\bibitem{qwen}
Qwen Team, Alibaba Group,
\emph{Qwen2.5 Technical Report},
arXiv:2412.15115, 2024.

\bibitem{ollama}
Ollama,
\emph{Run large language models locally},
\url{https://ollama.com/}.

\bibitem{k3s}
Rancher,
\emph{k3s --- Lightweight Kubernetes},
\url{https://k3s.io/}.

\bibitem{terraform}
HashiCorp,
\emph{Terraform documentation},
\url{https://developer.hashicorp.com/terraform/docs}.

\bibitem{ansible}
Red Hat,
\emph{Ansible documentation},
\url{https://docs.ansible.com/}.

\bibitem{tailscale}
Tailscale Inc.,
\emph{Tailscale documentation},
\url{https://tailscale.com/kb/}.

\end{thebibliography}

\end{document}
